\documentclass[11pt,a4paper]{article}

\usepackage[utf8]{inputenc}
\usepackage[T1]{fontenc}
\usepackage[brazil]{babel}
\usepackage[margin=2.4cm]{geometry}
\usepackage{amsmath,amssymb}
\usepackage{booktabs}
\usepackage{longtable}
\usepackage{xcolor}
\usepackage{graphicx}
\usepackage[hidelinks]{hyperref}

\newcommand{\code}[1]{\texttt{\small #1}}
\newcommand{\unit}[2]{\ensuremath{#1\,\mathrm{#2}}}

\title{\textbf{Avaliador de Vibração Humana}\\[0.3em]
\large Reimplementação e extensão do software de avaliação de risco\\
da ISO 2631-5:2018, com NHO-09, Diretiva 2002/44/CE e análise de sensibilidade\\[0.5em]
\normalsize Documentação técnica: cadeias de cálculo, engenharia reversa
dos dados,\\ validação, defeitos corrigidos, arquitetura atual e questões em aberto}
\author{GRAVI-SH --- Grupo de Acústica e Vibrações em Seres Humanos\\
Universidade Federal de Minas Gerais}
\date{\today}

\begin{document}
\maketitle

\begin{abstract}
\noindent
Este documento descreve o \textbf{Avaliador de Vibração Humana}, aplicativo do
GRAVI-SH que nasceu como reimplementação em Python do software de avaliação de
risco associado à ISO 2631-5:2018 --- exposição humana à vibração de corpo
inteiro contendo choques múltiplos --- e cresceu para um instrumento mais
abrangente de avaliação de exposição humana à vibração. O software original de
referência para a ISO 2631-5, distribuído como executável compilado a partir de
MATLAB, não estava operacional no ambiente de trabalho disponível na época da
reimplementação.

O núcleo cobre os \emph{dois} caminhos de avaliação previstos pela norma: o do
Anexo A com o Anexo E, para exposições sem perda de contato com o assento, e o
da Cláusula 5 com os Anexos C e D, para exposições com perda de contato. Inclui
ainda as sete ponderações em frequência da ISO 2631-1, cinco métodos de
estimação espectral para análise, e a verificação automática do critério de
validade que decide qual dos dois caminhos se aplica.

Em torno desse núcleo, o software cresceu para cobrir também: avaliação de
registros de campo do registrador do laboratório pela NHO-09 e pela Diretiva
2002/44/CE, com reprodução auditada do método do laboratório; um explorador de
sinais de propósito geral, independente da cadeia normativa; uma análise de
sensibilidade que varre os 30 modelos antropométricos e as sete posturas da
norma para quantificar a dispersão do resultado; e execução em segundo plano
para as operações pesadas, de modo que carregar dezenas de arquivos ou trocar de
postura não trava a interface.

O trabalho de reimplementação exigiu engenharia reversa do pacote de funções de
transferência distribuído com a norma, cuja documentação é incompleta, e a
identificação de uma etapa de processamento ausente do material disponível.
Treze casos de referência da norma são reproduzidos, e o método do laboratório
foi auditado e reproduzido exatamente a partir do seu código-fonte MATLAB. No
percurso foram encontrados e corrigidos defeitos que produziam erros de até
15\% no núcleo da norma e de até 5,7$\times$ no VDV do método do laboratório, e
permanecem questões que os dados disponíveis não resolvem por completo --- entre
elas a aplicabilidade do modelo biomecânico adulto a usuários pediátricos,
linha de pesquisa para a qual este documento já incorpora um levantamento de
literatura recente.
\end{abstract}

\tableofcontents
\newpage

% =====================================================================
\section{Contexto e objetivo}

A ISO 2631-5:2018 estabelece métodos para avaliar o risco à saúde da coluna
lombar decorrente da exposição ocupacional à vibração de corpo inteiro com
choques múltiplos. O método difere do da ISO 2631-1 por não se basear em
aceleração ponderada: estima, a partir da aceleração medida nos pontos de
contato entre o corpo e o veículo, as forças de compressão que atuam nos discos
intervertebrais lombares, e compara a dose acumulada ao longo da vida laboral
com a resistência da vértebra, que decai com a idade.

O objetivo original foi reconstruir esse software de forma operacional,
auditável e extensível, atendendo a:

\begin{itemize}
  \item cálculo de $S_d^A$ (Anexo A) e $R^A$ (Anexo E);
  \item cálculo de $D_z$, $S_d$, $R$ e da probabilidade de dano (Cláusula 5 com
        Anexo C);
  \item reprodução exata dos resultados de referência do software original;
  \item interface gráfica com visualização nos domínios do tempo e da
        frequência, com seleção individual de séries;
  \item transparência quanto ao processamento aplicado, permitindo a
        introdução de ponderações e filtros;
  \item empacotamento como executável autocontido.
\end{itemize}

O escopo cresceu desde então, e os objetivos abaixo se somaram aos originais:

\begin{itemize}
  \item avaliação de registros de campo (\code{.m}) pela NHO-09 da Fundacentro
        e pela Diretiva 2002/44/CE, com reprodução auditada do software de
        análise do laboratório (Seção~\ref{sec:campo});
  \item análise de sensibilidade postural e antropométrica, varrendo os 30
        modelos e as sete posturas do Anexo A num único sinal medido
        (Seção~\ref{sec:sensibilidade});
  \item explorador de sinais de propósito geral, tanto para arquivos ASCII de
        quatro colunas quanto para os registros \code{.m} do registrador de
        campo, independente da cadeia de cálculo normativa
        (Seções~\ref{sec:explorer} e~\ref{sec:fieldexplorer});
  \item execução em segundo plano para toda operação que envolva leitura ou
        recálculo de múltiplos arquivos, com barra de progresso, de modo que a
        interface permaneça responsiva (Seção~\ref{sec:threading});
  \item renomeação do produto: a ISO 2631-5 é hoje um entre vários caminhos de
        avaliação disponíveis no aplicativo, não o escopo inteiro dele --- daí
        o nome \emph{Avaliador de Vibração Humana}, adotado nesta versão.
\end{itemize}

O laboratório onde o trabalho se insere mede vibração de corpo inteiro em
usuários de cadeira de rodas, com medição triaxial em três pontos de contato:
assento, encosto e apoio de pés. Essa configuração mapeia diretamente sobre os
quatro pontos de contato do Anexo A, ficando o ponto correspondente às mãos sem
excitação.

% =====================================================================
\section{Fundamentação normativa}

\subsection{Os dois caminhos da norma e o critério que os separa}
\label{sec:validade}

A norma prevê dois métodos distintos, e a escolha entre eles não é do analista:
é determinada pela amplitude do sinal. A Seção 6.2 estabelece que, depois de
aplicados os filtros limitadores de banda da ISO 2631-1, os picos de aceleração
na direção $z$ não podem exceder \unit{9{,}81}{m/s^2} para que os Anexos A e E
se apliquem. Acima desse limite valem a Cláusula 5 e os Anexos C e D.

O limite corresponde a $1\,g$ e tem significado físico direto: acima dele
ocorre perda de contato entre a pessoa e o assento, e o modelo biodinâmico do
Anexo A, que pressupõe contato contínuo, deixa de valer.

Esta verificação estava ausente da versão original do software de referência.
Sua implementação revelou que \emph{dois dos arquivos de exemplo distribuídos
com a própria norma falham no teste}: o registro do bulldozer atinge
\unit{30{,}8}{m/s^2} ponderado e um dos registros de pá carregadeira também
excede o limite.

\subsection{Caminho do Anexo A com o Anexo E}

As forças intervertebrais resultam da superposição das contribuições de quatro
pontos de contato e três direções de excitação, no domínio da frequência. Para o
nível discal $\ell$, esta é a Fórmula (A.1):

\begin{equation}
F_\ell(f) \;=\; \sum_{p=1}^{4}\ \sum_{d=1}^{3}
      H_{\,p,d,\ell}(f)\; U_{p,d}(f)
\label{eq:a1}
\end{equation}

\noindent
onde $U_{p,d}(f)$ é o espectro de deslocamento da excitação no ponto $p$ e
direção $d$, e $H_{p,d,\ell}(f)$ a função de transferência complexa
correspondente.

A dose de compressão de uma exposição, Fórmula (A.2):

\begin{equation}
S^{A}_\ell \;=\; \left[\ \sum_{i}
    \left(\frac{C_{\mathrm{dyn},i}}{B_\ell}\right)^{6}\ \right]^{1/6}
\label{eq:a2}
\end{equation}

\noindent
Os picos $C_{\mathrm{dyn},i}$ são o valor extremo da força adicional de
compressão entre dois cruzamentos consecutivos do valor médio, e são negativos.
$B_\ell$ é a área média de placa terminal do nível. O expoente 6 corresponde à
regra de acúmulo de dano de Palmgren--Miner.

A dose diária equivalente, a partir de $j$ exposições medidas, Fórmula (A.3):

\begin{equation}
S^{A}_{d,\ell} \;=\; \left[\ \sum_{j}
    \left(S^{A}_{\ell,j}\right)^{6}\, \frac{t_{d,j}}{t_{m,j}}\ \right]^{1/6}
\label{eq:a3}
\end{equation}

\noindent
$t_{m,j}$ é a duração da medição e $t_{d,j}$ o tempo diário de exposição naquela
condição. A razão $t_d/t_m$ é um fator de extrapolação: assume-se que o trecho
medido representa o dia inteiro. Como o expoente é 6, a extrapolação é suave ---
multiplicar a exposição por 36 multiplica a dose por $36^{1/6} \approx 1{,}82$.

Para padrões de exposição que variam ao longo do ano, a Fórmula (E.3) pondera
pelos dias de cada padrão:

\begin{equation}
S^{A}_{q,\ell} \;=\; \left[\ \sum_{h}
    \left(S^{A}_{d,\ell,h}\right)^{6}\, \frac{N_h}{N}\ \right]^{1/6},
\qquad N = \sum_h N_h
\label{eq:e3}
\end{equation}

O fator de risco acumula a dose ao longo dos anos, Fórmula (E.1):

\begin{equation}
R^{A}_\ell \;=\; \left[\ \sum_{i=1}^{n}
   \left( \frac{S^{A}_{d,\ell,i}\; N_i^{1/6}}
               {S^{A}_{u,i} - S^{A}_{\mathrm{stat},\ell,i}} \right)^{6}
   \ \right]^{1/6}
\label{eq:e1}
\end{equation}

\noindent
com a resistência última dada pela Fórmula (E.2):

\begin{equation}
S^{A}_{u,i} \;=\; 6{,}765024\ \mathrm{MPa}
   \;-\; 0{,}067184\ \mathrm{MPa}\cdot (b+i)
\label{eq:e2}
\end{equation}

\noindent
onde $b$ é a idade no início da exposição e $i$ percorre os anos. A relação foi
derivada de estudos \emph{in vitro} e considera a proteção de 50\% da população
de motoristas. A norma registra explicitamente que a Fórmula (E.2) e a Fórmula
(C.4) do Anexo C foram derivadas por caminhos diferentes e não são idênticas.

A interpretação segue os limiares do Anexo E: $R^A < 0{,}8$ indica baixa
probabilidade de efeito adverso, $R^A > 1{,}2$ indica alta. A faixa intermediária
fica sem veredito, o que é a forma adotada pela norma para acomodar a
variabilidade interindividual.

\subsection{Caminho da Cláusula 5 com o Anexo C}

Aqui a resposta da coluna é modelada por uma única função de transferência
assento-coluna, com um zero complexo e seis polos complexos, Fórmula (1):

\begin{equation}
H(\omega) \;=\;
\frac{1 + 2\zeta_1 \dfrac{s}{\omega_1} + \left(\dfrac{s}{\omega_1}\right)^2}
     {\displaystyle\prod_{n=2}^{7}
      \left[1 + 2\zeta_n \dfrac{s}{\omega_n} + \left(\dfrac{s}{\omega_n}\right)^2\right]},
\qquad s = j\omega
\label{eq:c5tf}
\end{equation}

\noindent
com os parâmetros da Tabela~\ref{tab:poles}. A norma registra que a
transmissibilidade em corrente contínua foi ajustada para 1,0; a implementação
devolve exatamente 1,000000, o que serve de verificação.

\begin{table}[h]
\centering
\caption{Parâmetros da Fórmula (1), em rad/s.}
\label{tab:poles}
\begin{tabular}{lcccccccc}
\toprule
 & $\omega_1$ & $\omega_2$ & $\omega_3$ & $\omega_4$ & $\omega_5$ & $\omega_6$ & $\omega_7$ \\
\midrule
$\omega$ & 34 & 31 & 230 & 260 & 320 & 380 & 420 \\
$\zeta$  & 0,35 & 0,21 & 0,88 & 0,80 & 0,40 & 0,75 & 0,65 \\
\bottomrule
\end{tabular}
\end{table}

A dose de aceleração, Fórmula (3), conta apenas os picos \emph{positivos} entre
cruzamentos por zero:

\begin{equation}
D_z \;=\; 1{,}07 \left[\ \sum_i A_{z,i}^{6}\ \right]^{1/6}
\label{eq:dz}
\end{equation}

A dose diária normaliza pelo tempo de exposição, Fórmula (5), e a dose de
compressão sai da Fórmula (C.1) com o fator $m_z$ da Fórmula (C.2):

\begin{equation}
S_d = m_z\,D_{zd},
\qquad
m_z = \frac{\text{massa corporal} \times \text{fração acima da pelve}}
           {\text{área de placa terminal}}
\label{eq:mz}
\end{equation}

A variável de tensão, Fórmula (C.3), e a resistência última, Fórmula (C.4):

\begin{equation}
R = \left[\ \sum_{i=0}^{n-1}
    \left(\frac{S_{d,i}\,N_i^{1/6}}{S_{u,i} - S_{\mathrm{stat},i}}\right)^{6}
    \ \right]^{1/6},
\qquad
S_{u,i} = 6{,}75\ \mathrm{MPa} - S_{\mathrm{age}}\,(b+i)
\label{eq:c3}
\end{equation}

\noindent
com $S_{\mathrm{age}}$ de \unit{0{,}052}{MPa} para homens e \unit{0{,}039}{MPa}
para mulheres. A probabilidade de dano segue um modelo de Weibull, Fórmula
(C.5):

\begin{equation}
\Pi(R) \;=\; 1 - \exp\!\left[-\left(\frac{R}{\alpha}\right)^{\beta}\right]
\label{eq:weibull}
\end{equation}

\noindent
com $\alpha = 1{,}613$ e $\beta = 2{,}799$ para homens, $\alpha = 0{,}959$ e
$\beta = 3{,}709$ para mulheres. Os limiares de interpretação aqui são
$\Pi < 10\%$ e $\Pi > 50\%$.

Uma diferença de escopo merece registro: este caminho usa \emph{apenas} a
direção $z$ medida no assento. Encosto, pés e as direções horizontais não
entram. Em compensação, o $m_z$ é parâmetro aberto, o que torna o método
adaptável a corpos fora do padrão de operador adulto --- ao contrário do Anexo
A, cujos modelos antropométricos são fixos. Esse detalhe volta a importar na
discussão sobre aplicação a crianças (Seção~\ref{sec:criancas}).

% =====================================================================
\section{Cadeia de cálculo do Anexo A}
\label{sec:cadeiaA}

Cada etapa corresponde a uma função identificável em \code{iso26315/core.py}.

\subsection{Leitura do registro}

Arquivos ASCII de quatro colunas sem cabeçalho: tempo em segundos, seguido de
$a_x$, $a_y$ e $a_z$ em \unit{}{m/s^2}, no sistema basicêntrico da ISO 2631-1.
O passo de tempo é inferido do primeiro e do último instante, e a uniformidade
da base de tempo é verificada. Registros de comprimento ímpar perdem a última
amostra, pois a formulação do espectro unilateral exige comprimento par.

\subsection{Transformada de Fourier: definição exata}
\label{sec:fft}

A transformada é aplicada ao registro inteiro de uma só vez, com janela
retangular, sem segmentação, sobreposição, preenchimento com zeros ou média de
blocos. Para $N$ amostras (par), passo $\Delta t$ e $M = N/2$:

\begin{equation}
A_k \;=\; \frac{1}{M}\sum_{n=0}^{N-1} a_n\, e^{-2\pi i kn/N},
\qquad k = 0,\dots,M
\label{eq:fft}
\end{equation}

\noindent
com $A_0$ e $A_M$ multiplicados por $1/2$. Essa normalização produz um espectro
unilateral de \emph{amplitude}: uma senoide pura de amplitude $a$ resulta em um
bin de módulo $a$. A grade é $f_k = k/(N\Delta t)$, com resolução
$\Delta f = 1/(N\Delta t)$ e Nyquist $1/(2\Delta t)$. A transformada inversa
reverte exatamente essa normalização, reconstruindo o sinal real por simetria
hermitiana.

\subsection{Limitação de banda}

Os bins abaixo de \unit{0{,}5}{Hz} são zerados, conforme a Cláusula 5 da norma,
eliminando deriva e componente contínua. Os bins acima de \unit{100}{Hz} são
descartados. Como a reconstrução no tempo parte do espectro truncado, a taxa de
amostragem efetiva do sinal de força reconstruído passa a ser \unit{200}{Hz}.

\subsection{Ausência de ponderação em frequência}

Nenhuma ponderação é aplicada nesta cadeia. O Anexo A trabalha com aceleração
não ponderada: os filtros $W_k$ e $W_d$ pertencem à ISO 2631-1 e não fazem parte
deste método, porque a dependência em frequência da resposta biomecânica já está
incorporada às funções de transferência do modelo de elementos finitos.

\subsection{Conversão para deslocamento}
\label{sec:disp}

\textbf{Esta é a etapa crítica do trabalho e a única ausente do material
disponível.} As funções de transferência do Anexo A mapeiam \emph{deslocamento}
em força, não aceleração em força. O espectro medido deve ser integrado duas
vezes no domínio da frequência:

\begin{equation}
U(f) \;=\; \frac{A(f)}{-(2\pi f)^{2}}
\label{eq:disp}
\end{equation}

A omissão produz forças cerca de mil vezes maiores. A
Seção~\ref{sec:descoberta} descreve como a etapa foi identificada.

\subsection{Interpolação das funções de transferência}
\label{sec:fase}

As funções de transferência do modelo são definidas em grade uniforme de
\unit{0{,}25}{Hz} a \unit{30}{Hz}, com 120 pontos. A grade da excitação é
determinada pelo comprimento do registro e é, em geral, muito mais fina. A
interpolação é por \emph{spline} cúbica, separadamente sobre o módulo e sobre a
fase.

Fora da faixa do modelo o resultado é identicamente nulo. Não há extrapolação:
abaixo de \unit{0{,}25}{Hz} e acima de \unit{30}{Hz} o modelo não tem validade.

A interpolação da fase é numericamente frágil, e o assunto é tratado em detalhe
na Seção~\ref{sec:fasebug}.

\subsection{Superposição, retorno ao tempo e contagem de picos}

Aplica-se a Equação~\eqref{eq:a1}. Das três componentes de força resultantes, a
terceira é a compressão axial e é a usada no cálculo da dose.

Após a transformada inversa, o valor médio do sinal de força é substituído pela
força estática do modelo de elementos finitos, tabelada por postura,
antropometria e nível discal. Essa força estática também define
$S^{A}_{\mathrm{stat}}$.

Os picos são identificados entre cruzamentos consecutivos do valor médio: o
sinal é centrado, determinam-se os índices de troca de sinal, e em cada
intervalo cuja primeira amostra seja negativa registra-se o mínimo.

% =====================================================================
\section{Cadeia de cálculo da Cláusula 5}
\label{sec:cadeiaC5}

As duas cadeias são implementadas em módulos independentes e não compartilham
condicionamento de sinal, porque os requisitos conflitam
(Tabela~\ref{tab:cadeias}).

\begin{table}[h]
\centering
\caption{Diferenças entre as duas cadeias.}
\label{tab:cadeias}
\begin{tabular}{lll}
\toprule
 & \textbf{Cláusula 5 / Anexo C} & \textbf{Anexo A / Anexo E} \\
\midrule
passa-alta            & \unit{0{,}1}{Hz}, 2 polos & \unit{0{,}5}{Hz} \\
passa-baixa           & \unit{80}{Hz}, 4 polos    & \unit{100}{Hz} \\
filtros da ISO 2631-1 & proibidos                 & só na validação \\
direções              & só $z$, no assento        & $x$, $y$, $z$ em 4 pontos \\
excitação             & aceleração                & deslocamento \\
saída                 & $D_z$, $S_d$, $R$, $\Pi$  & $S_d^A$, $R^A$ \\
\bottomrule
\end{tabular}
\end{table}

O corte inferior em \unit{0{,}1}{Hz} existe para preservar a queda livre antes
do impacto: filtrar em \unit{0{,}5}{Hz} empurraria o $-1\,g$ de volta a zero
antes do choque e erraria o pico em até $1\,g$. Os filtros são causais,
coerentes com a implementação digital em equação de diferenças do Anexo D.

A resposta da coluna é obtida aplicando a Equação~\eqref{eq:c5tf} ao espectro do
sinal condicionado e transformando de volta. Os picos positivos entre
cruzamentos por zero alimentam a Equação~\eqref{eq:dz}.

\subsection{Detecção de perda de contato}
\label{sec:contato}

A norma exige detectar os trechos sem contato e excluí-los da exposição,
sugerindo chave de contato, vídeo ou método equivalente. O arranjo experimental
disponível não dispõe desse instrumento.

O software oferece uma \textbf{heurística que não consta da norma}: marca como
queda livre os trechos em que a aceleração vertical permanece abaixo de
\unit{-7}{m/s^2} por pelo menos \unit{0{,}15}{s}. A opção pode ser desligada, os
limiares são parâmetros, e a interface informa que os trechos marcados devem ser
conferidos antes de aceitos. Qualquer resultado obtido com a heurística ligada
deve registrar esse fato.

% =====================================================================
\section{Ponderações em frequência da ISO 2631-1}

Sete ponderações estão implementadas a partir das Tabelas A.1 e A.2 daquela
norma: $W_k$, $W_d$, $W_c$, $W_e$, $W_j$, $W_f$ e $W_b$. Cada uma é o produto
dos quatro fatores da Seção A.2 --- passa-alta e passa-baixa Butterworth de dois
polos, transição aceleração-velocidade e degrau ascendente --- montado no
domínio de Laplace:

\begin{align}
H_h(s) &= \frac{1}{1 + \sqrt{2}\,\omega_1/s + (\omega_1/s)^2} \\
H_l(s) &= \frac{1}{1 + \sqrt{2}\,s/\omega_2 + (s/\omega_2)^2} \\
H_t(s) &= \frac{1 + s/\omega_3}{1 + s/(Q_4\omega_4) + (s/\omega_4)^2} \\
H_s(s) &= \frac{1 + s/(Q_5\omega_5) + (s/\omega_5)^2}
               {1 + s/(Q_6\omega_6) + (s/\omega_6)^2}
          \left(\frac{\omega_5}{\omega_6}\right)^2
\end{align}

\subsection{Validação e uma inconsistência na tabela da norma}

A validação usa a Tabela 3 da ISO 2631-1, que dá os fatores em bandas de um
terço de oitava. O erro relativo mediano é de 0,05\% para o $W_k$ e 0,04\% para
o $W_d$, compatível com o arredondamento de três dígitos da tabela.

Há um ponto discrepante que merece registro. No $W_d$ a \unit{10}{Hz} a
divergência sobe para 4,9\%, mas a causa não é a implementação: a coluna de
fatores da tabela traz 212, enquanto a coluna em decibéis da mesma linha traz
$-13{,}91$\,dB, que corresponde a 201,6. Os dois valores da própria tabela são
mutuamente inconsistentes, e o cálculo coincide com a coluna em decibéis, dando
201,7. Muito provavelmente trata-se de erro do reconhecimento óptico usado para
extrair a norma; convém conferir no documento original.

A ponderação $W_c$, usada na avaliação de conforto do encosto
(Seção~\ref{sec:posicao}), foi validada contra a Tabela 4 da ISO 2631-1, com erro
mediano de 0,02\%.

\subsection{Aplicação e escopo}

\textbf{As ponderações não entram em nenhum dos dois cálculos de risco do
Anexo A/Cláusula 5.} A única aplicação normativa nessa cadeia é dentro da
verificação do critério de \unit{9{,}81}{m/s^2}, onde o $W_k$ é obrigatório.
Fora disso, servem para análise comparativa com as métricas da ISO 2631-1 no
explorador de sinais e no painel de campo (Seção~\ref{sec:campo}), e a interface
repete esse aviso ao lado do seletor.

O painel de ponderação permite escolher filtro e fator $k$ por eixo, com sete
predefinições: sem ponderação (o que a ISO 2631-5 exige), saúde da ISO 2631-1
($W_d/W_d/W_k$ com $k = 1{,}4/1{,}4/1{,}0$), conforto no assento, conforto no
encosto, enjoo de movimento ($W_f$ só em $z$) e conforto da BS 6841. Nas
predefinições os controles ficam travados; em modo \emph{custom} tanto o filtro
quanto o $k$ são livres, e o rodapé do painel avisa que os valores deixam de
corresponder a qualquer norma. O painel aparece tanto na aba de domínio da
frequência quanto no explorador de sinais.

% =====================================================================
\section{Métodos de estimação espectral}
\label{sec:espectral}

O espectro de bloco único da Seção~\ref{sec:fft} é o que alimenta o cálculo,
mas é ruidoso para análise visual: um periodograma de realização única tem
variância comparável à própria média, independentemente do comprimento do
registro. Cinco métodos estão disponíveis para visualização.

\begin{table}[h]
\centering
\caption{Métodos de estimação espectral disponíveis.}
\begin{tabular}{ll}
\toprule
\textbf{Método} & \textbf{Uso} \\
\midrule
Bloco único & FFT complexa do registro inteiro; alimenta o cálculo \\
Welch, amplitude & Média de periodogramas, corrigida em amplitude \\
Welch, PSD & Densidade espectral de potência, em $(\mathrm{m/s^2})^2/\mathrm{Hz}$ \\
Welch, RMS por banda & PSD integrada em bandas de largura escolhida \\
Um terço de oitava & Bandas com centros nominais da ISO 266 \\
\bottomrule
\end{tabular}
\end{table}

Janelas disponíveis: retangular, Hann, Hamming, Flat top e Blackman-Harris. O
comprimento de segmento e a sobreposição são configuráveis.

\subsection{Calibração verificada}

Uma senoide de amplitude 1,0 a \unit{5}{Hz} devolve pico de 1,0000 com a janela
Flat top e 0,9959 com Hann --- coerente com o propósito da Flat top, que é
precisão de amplitude à custa de resolução em frequência. A densidade espectral
de ruído branco de variância unitária amostrado a \unit{500}{Hz} devolve
0,00399 contra o valor teórico $2\Delta t = 0{,}00400$.

\subsection{Por que o autoespectro não pode alimentar o cálculo}

A média de Welch reduz a variância do estimador promediando módulos ao quadrado,
operação que \emph{descarta a fase}. Sem fase não existe transformada inversa, e
portanto não há como reconstruir a força no domínio do tempo --- que é onde os
picos de compressão das Equações~\eqref{eq:a2} e \eqref{eq:dz} são contados.

Janelar o registro inteiro antes da FFT única também não serve: a janela atenua
as extremidades, e picos de choque que caiam nos primeiros ou últimos segundos
seriam artificialmente reduzidos, justamente os eventos que dominam uma soma de
sexta potência.

Existe um caminho legítimo para janelar sem perder a reconstrução ---
processamento segmentado com sobreposição e soma, respeitando a condição COLA
--- mas ele alteraria os resultados em relação à referência validada e não foi
implementado.

\subsection{Fase e coerência entre pontos de contato}
\label{sec:fasecoerencia}

Dois recursos de diagnóstico foram adicionados nesta versão, ambos de
visualização apenas --- nenhum dos dois entra no cálculo de $S_d^A$ ou $R^A$,
pela mesma razão da Seção~\ref{sec:espectral}: o método já usa a fase completa
do espectro de bloco único internamente, e esses recursos só a expõem ao
usuário.

\paragraph{Fase (canal único).} A opção \emph{Excitation phase spectrum}, no
mesmo seletor de gráfico da aba \emph{Frequency domain}, mostra
$\arg\big(A(f)\big)$ desenrolado, do espectro complexo de bloco único que
alimenta a Fórmula (A.1) --- o único espectro do software que preserva fase.
É por isso que esta visualização não aceita os métodos de Welch/oitava: a
média de Welch descarta fase (Seção~\ref{sec:espectral}), então não há uma
``fase de Welch'' coerente a mostrar. Disponível também no Signal explorer
(registros ASCII).

\paragraph{Coerência (par de canais).} A opção \emph{Coherence between
contact points} calcula a coerência magnitude-quadrado
\begin{equation}
\gamma^2_{xy}(f) \;=\; \frac{|S_{xy}(f)|^2}{S_{xx}(f)\,S_{yy}(f)} \in [0, 1]
\label{eq:coherence}
\end{equation}
entre um ponto de contato de referência e cada um dos outros três, na mesma
grade de frequência, com $S_{xy}$, $S_{xx}$ e $S_{yy}$ estimados por médias de
segmentos de Welch (o mesmo mecanismo de \code{welch\_spectrum}, generalizado
para densidade espectral cruzada em \code{core.welch\_cross\_spectrum}). Ao
contrário da fase de canal único, a coerência --- e a fase cruzada que a
acompanha, $\arg(S_{xy})$ --- só é informativa com mais de um segmento
promediado: com um único periodograma, $\gamma^2 = 1$ em todo bin por
construção, o que nada diz sobre a relação real entre os sinais. A interface
avisa quando o número de segmentos médios é baixo (Segment length curto
produz mais segmentos e uma estimativa mais confiável, à custa de resolução em
frequência).

\begin{quote}
\textbf{Por que só entre os quatro pontos de uma mesma exposição.} A
Equação~\eqref{eq:coherence} só tem interpretação física quando $x$ e $y$
compartilham grade de tempo comum --- exatamente a condição que
\code{core.check\_time\_base} já verifica para poder somar os quatro pontos na
Fórmula (A.1) (Seção~9.5). É por isso que o seletor de ``ponto de referência''
trabalha dentro de uma única exposição carregada, e não entre arquivos
avulsos. O Signal explorer generalista também oferece coerência (contra o
primeiro arquivo marcado), mas, ao contrário da aba \emph{Frequency domain},
não tem como verificar se os arquivos carregados ali foram de fato gravados em
conjunto --- a nota exibida junto ao gráfico avisa disso explicitamente. Para
os registros \code{.m} do registrador de campo, coerência entre posições
\emph{não} foi implementada: a Seção~\ref{sec:conjuntoB} já documenta que os
relógios das diferentes unidades não estão sincronizados, o que tornaria
qualquer coerência entre posições um artefato do desalinhamento temporal, não
uma medida real de correlação mecânica.
\end{quote}

% =====================================================================
\section{Engenharia reversa do pacote de dados}
\label{sec:dados}

O pacote distribuído com a norma contém seis arquivos compactados, um por grupo
de postura, cada um com 810 arquivos de texto, além de um arquivo de forças
estáticas e um de descrição da estrutura. A descrição informa apenas que a
matriz principal é um arranjo tridimensional de células $30\times4\times3$, e
que cada célula contém $3\times121\times7$ entradas do tipo \emph{double}. Não
informa o significado das dimensões, nem que os valores são complexos.

\subsection{Estrutura decodificada}

Cada arquivo \code{BSKUefuM\_k\_p\_d} é um CSV de 3 linhas por 847 colunas:

\begin{center}
\begin{tabular}{ll}
\toprule
\textbf{Dimensão} & \textbf{Significado} \\
\midrule
847 colunas & 7 blocos de 121 valores \\
\quad bloco 1 & vetor de frequência, \unit{0{,}25}{} a \unit{30}{Hz} \\
\quad blocos 2 a 7 & os seis níveis discais, T12/L1 a L5/S1 \\
\quad coluna 1 de cada bloco & força estática do modelo de elementos finitos \\
\quad colunas 2 a 121 & função de transferência complexa \\
3 linhas & direção da \emph{força}: $x$, $y$, compressão \\
índice $p$ (1 a 4) & ponto de excitação: assento, pés, mãos, encosto \\
índice $d$ (1 a 3) & direção da \emph{excitação}: $x$, $y$, $z$ \\
índice $k$ (1 a 30) & modelo antropométrico (ver Seção~\ref{sec:indice}) \\
\bottomrule
\end{tabular}
\end{center}

A atribuição das linhas à direção da \emph{força}, e não da excitação, foi
estabelecida por argumento de simetria: no arquivo correspondente à excitação em
$y$, apenas a segunda linha apresenta magnitude apreciável, sendo as demais
praticamente nulas. Excitação lateral gera força lateral, não compressão.

A natureza da primeira coluna de cada bloco foi confirmada no código-fonte
original, que a atribui diretamente aos resultados estáticos do modelo.

O arquivo de forças estáticas contém 420 linhas, correspondentes a 7 posturas
$\times$ 10 antropometrias $\times$ 6 níveis discais, com três colunas de força.
O índice de acesso é $60\,(\mathrm{postura}-1) + 6\,(\mathrm{antropometria}-1)$.

\subsection{Um sexto modelo de postura fora da tabela da norma}

A Tabela A.1 da norma define cinco grupos de postura, a partir de sete ângulos
corporais medidos em operadores de empilhadeira, pá carregadeira, escavadeira,
forwarder e harvester. O pacote de dados, porém, traz \emph{seis} conjuntos, e o
sexto se chama \code{handsinlap} --- mãos no colo. Ele não consta da tabela
normativa, mas o código original o trata como postura 7. O software conta,
portanto, com sete posturas ao todo --- a informação que alimenta a análise de
sensibilidade postural (Seção~\ref{sec:sensibilidade}).

A inspeção dos dados confirma o que o nome sugere: nesse modelo o ganho do
ponto de contato correspondente ao encosto é exatamente zero, e o ganho do
assento é cerca de 40\% menor que nos demais. É um modelo de pessoa sentada sem
apoio de costas e sem as mãos em qualquer comando.

Para avaliação de usuários de cadeira de rodas, nenhum dos sete modelos é
adequado: todos pressupõem pés apoiados no piso da cabine ou em pedais, com o
quadril mais alto que os joelhos. O modelo \code{handsinlap} acerta a condição
das mãos mas descarta a medição do encosto. A recomendação metodológica é
executar os sete e reportar a dispersão --- a força estática de compressão em
L3/L4 varia 67\% entre a postura de menor e a de maior carga para a mesma
pessoa, o que dá a escala da incerteza envolvida. Essa recomendação passou a ser
uma funcionalidade do software (Seção~\ref{sec:sensibilidade}), e não apenas uma
prática manual.

\subsection{Identificação da conversão para deslocamento}
\label{sec:descoberta}

A primeira implementação, tratando as funções de transferência como
aceleração\,$\rightarrow$\,força, produziu $S_d^A$ da ordem de
\unit{500}{MPa}, contra os \unit{0{,}58}{MPa} de referência --- três ordens de
grandeza de divergência.

A causa foi identificada por inspeção dimensional. Para o modelo antropométrico
do caso de referência, a função assento\,$z$\,$\rightarrow$\,compressão em L4/L5
apresenta módulo de aproximadamente $4{,}6\times10^{4}$ em \unit{5{,}25}{Hz}.
Interpretada como força por aceleração, essa grandeza corresponderia a uma massa
aparente de \unit{46}{t}, fisicamente impossível. Dividindo por $\omega^2$,
obtém-se \unit{42{,}6}{kg}, compatível com a massa efetiva do tronco de um
indivíduo de \unit{110}{kg}. A verificação repetida em várias frequências
resultou sempre em massas efetivas entre \unit{10}{} e \unit{60}{kg}.

Incorporada a Equação~\eqref{eq:disp}, os resultados passaram a coincidir com a
referência em todos os dígitos publicados. A hipótese é ainda corroborada por um
comentário no código-fonte original, que se refere à excitação como
deslocamento. Note-se que a etapa está ausente do trecho de código-fonte
disponível: a extração por reconhecimento óptico está degradada em várias
passagens, e a função de importação dos registros é uma delas.

\subsection{Conversão para formato compacto}

Os 810 arquivos de texto de cada postura, totalizando cerca de 172\,MB, foram
convertidos para arranjos binários comprimidos de precisão simples complexa,
totalizando 35\,MB. A conversão é feita uma única vez por um utilitário incluído
no projeto (\code{tools/build\_data.py}), e a validação foi repetida depois dela
sem alteração nos resultados.

% =====================================================================
\section{Validação}

\subsection{Anexo A com Anexo E: oito casos}

\begin{center}
\begin{tabular}{llll}
\toprule
\textbf{Caso} & \textbf{Massa} & \textbf{Erro relativo} & \textbf{Situação} \\
\midrule
Orientadora (pá carregadeira) & 110\,kg & $3{,}9\times10^{-5}$ & exato \\
Individual (bulldozer)        & 110\,kg & $2{,}2\times10^{-5}$ & exato \\
Individual 1 a 4              & 110\,kg & $7\times10^{-6}$ a $3{,}3\times10^{-5}$ & exato \\
Individual 5 e 6              & 70 a 110\,kg & $5{,}6$ a $6{,}4\times10^{-4}$ & aproximado \\
\bottomrule
\end{tabular}
\end{center}

O caso orientador, detalhado:

\begin{center}
\begin{tabular}{lcccccc}
\toprule
$R^A$ & T12/L1 & L1/L2 & L2/L3 & L3/L4 & L4/L5 & L5/S1 \\
\midrule
Implementação & 0,98423 & 1,08503 & 1,22263 & 1,31024 & 1,28424 & 1,16975 \\
Referência    & 0,98423 & 1,085   & 1,2226  & 1,3102  & 1,2842  & 1,1698  \\
\bottomrule
\end{tabular}
\end{center}

\noindent
A dose no nível crítico resultou em $S_d^A = \unit{0{,}57776}{MPa}$, coincidente
com a referência.

Os seis casos com massa constante de \unit{110}{kg} fecham no nível do
arredondamento dos valores publicados. Os dois que variam a massa ficam uma
ordem de grandeza acima, pelo motivo discutido na Seção~\ref{sec:indice}. Esta
validação é reproduzível a qualquer momento por \code{python
tests/test\_reference.py}, e roda como parte da rotina de desenvolvimento antes
de qualquer entrega.

\subsection{Cláusula 5 com Anexo C: cinco verificações}

O exemplo numérico do fim do Anexo C --- homem de \unit{82}{kg},
\unit{120}{dias/ano}, dos 20 aos 40 anos, cinco acelerações de \unit{40}{m/s^2}
por dia --- reproduz $S_d = \unit{1{,}623}{MPa}$ e $R = 1{,}22$ exatamente,
usando o $m_z$ arredondado de 0,029 que a própria norma emprega. Com o $m_z$
exato de 0,02870 sai 1,606 e 1,207; a diferença é integralmente o arredondamento
da norma, não da implementação.

A Tabela C.2 reproduz exatamente nos dois sexos: $R$ de 0,72 / 1,42 / 2,17 para
10\%, 50\% e 90\% de risco em homens, e 0,52 / 0,87 / 1,20 em mulheres. O ganho
em corrente contínua da Fórmula (1) dá 1,000000.

\subsection{Ponderações}

Erro relativo mediano de 0,05\% ($W_k$) e 0,04\% ($W_d$) contra a Tabela 3 da
ISO 2631-1, com a ressalva da inconsistência da própria tabela a
\unit{10}{Hz}; e 0,02\% ($W_c$) contra a Tabela 4.

\subsection{Método do laboratório: dois casos}

Além dos treze casos normativos, o painel que reproduz o aplicativo
\code{WBV\_HSE\_v16} (Seção~\ref{sec:posicao}) foi validado contra dois casos
externos: a tela do próprio aplicativo, em 20/01/2025 às 8\,h ($A(8) =
0{,}303$), e o resultado de encosto publicado no artigo do CONEM2026
($A(8) = 0{,}685$, $VDV = 165{,}9$). Ambos são reproduzidos exatamente pelo
software, incluindo os defeitos do aplicativo original quando a caixa
\emph{Reproduce the app VDV defects} está marcada (Seção~\ref{sec:auditoria}).

% =====================================================================
\section{Defeitos encontrados e corrigidos}

Esta seção registra os erros identificados no percurso, por ordem de gravidade.
Todos foram corrigidos; os números indicam o que estava em jogo.

\subsection{Anos de afastamento descartados da linha do tempo}

\textbf{Erro de até 15\%.} Anos sem padrão de exposição estavam sendo filtrados
antes de montar o somatório da Equação~\eqref{eq:e1}. O índice $i$ desse
somatório é a idade, e a resistência última decai com ela; pular um ano vazio
comprimia a carreira inteira e fazia a pessoa envelhecer menos do que deveria,
\emph{subestimando} o risco.

O defeito só apareceu quando o caso de exemplo \code{indiv6}, que contém quatro
anos de afastamento, ficou disponível. Com os anos vazios contando a idade mas
não acumulando dose, o erro caiu de 15\% para $6{,}4\times10^{-4}$. O defeito
atingiria qualquer histórico com licença, troca de função ou desemprego.

\subsection{Soma de amplitudes entre eixos}

\textbf{Superestimativa de 51\%.} Na opção de somar os três eixos nos métodos de
Welch, as amplitudes estavam sendo somadas escalarmente. Amplitudes somam em
quadratura; potências é que somam escalarmente. Medido em um registro real:
0,677 contra 0,448 corretos.

Defeito relacionado: na análise de um terço de oitava com todos os eixos, os
sinais crus $a_x + a_y + a_z$ estavam sendo somados amostra a amostra antes da
transformada, o que permite cancelamento de fase entre eixos e não tem sentido
físico. Corrigido para somar as potências de banda.

Ambos afetavam apenas gráficos, nunca o cálculo de risco.

\subsection{Esgotamento de memória em históricos longos}

Um resultado completo guarda a série de força e o espectro de cada nível discal,
somando dezenas de MB por combinação antropometria $\times$ exposição. O caso
\code{indiv6}, com treze mudanças de massa corporal, derrubava o processo. O
mesmo padrão de crescimento voltou a aparecer, em escala maior, na análise de
sensibilidade (Seção~\ref{sec:sensibilidade}), que multiplica esse produto
cartesiano por 7 posturas $\times$ 30 antropometrias.

O cache passou a guardar o resultado completo apenas para as quatro primeiras
combinações --- as que alimentam os gráficos --- e um resumo com $S$,
$S_{\mathrm{stat}}$ e $t_m$ para as demais, que é tudo de que a dose e o risco
precisam. O espectro de deslocamento é liberado logo após o cálculo das forças.
O pico caiu para 2,1\,GB, com resultado inalterado. A análise de sensibilidade
reaproveita o mesmo cache, o que a mantém viável em memória apesar de rodar 210
combinações por sinal.

\subsection{Ausência do critério de validade}

A verificação dos \unit{9{,}81}{m/s^2} simplesmente não existia, de modo que o
software calculava e exibia resultados fora da faixa de validade do método sem
qualquer sinalização. Corrigido: a verificação é automática e o aviso aparece
acima da tabela de resultados.

\subsection{Consistência da base de tempo entre sensores}

A Equação~\eqref{eq:a1} soma as contribuições espectro a espectro, o que exige
grade de frequência comum. Sensores independentes derivam entre si. O software
passou a verificar antes de calcular: taxas de amostragem diferentes são
recusadas; deriva acima de 1\% é recusada com as durações de cada arquivo na
mensagem; deriva abaixo de 1\% trunca ao menor comprimento comum com aviso
visível informando quantas amostras foram descartadas. A tolerância está em
\code{core.TIME\_BASE\_TOLERANCE}.

\subsection{Artefatos numéricos da transformada: estudo controlado}
\label{sec:leakage}

A cadeia do Anexo A transforma, multiplica pela função de transferência e
transforma de volta. Nessa estrutura o \emph{leakage} clássico --- espalhamento
de uma senoide para bins vizinhos --- praticamente não importa, porque bins
individuais nunca são lidos: o espectro é apenas intermediário, e o
espalhamento se desfaz na volta. O truncamento do espectro, porém, introduz dois
artefatos reais, medidos aqui em sinais sintéticos de resposta conhecida.

\paragraph{Convolução circular.} Multiplicar no domínio da frequência equivale a
convolução circular no tempo: sem acolchoamento, a resposta ao impulso do fim do
registro enrola para o começo. Com um choque posicionado no último terço de
segundo de um registro de 30\,s, 30\% da amplitude de pico aparece nos primeiros
\unit{0{,}5}{s} do sinal reconstruído --- força que nunca existiu ali. Acolchoar
o registro com zeros reduz isso a 0,74\%.

\paragraph{Oscilação de Gibbs.} Zerar tudo abaixo de \unit{0{,}5}{Hz} e
descartar tudo acima de \unit{100}{Hz} é um filtro de parede vertical, cuja
resposta no tempo é um seno cardinal, não causal e oscilante. Um choque isolado
de \unit{5}{m/s^2} passa a produzir 16 picos contados onde existia um.

Aqui o estudo contrariou a expectativa inicial. Substituir o corte vertical por
uma transição de cosseno elevado \textbf{não reduz a contagem de picos}: os 16
permanecem, e a dose cai apenas 3,2\% no caso mais adverso, 0,4\% no trem de
choques e menos de 0,1\% quando os choques estão imersos em vibração contínua.
Ou seja, a multiplicação de picos é propriedade da limitação de banda em si, não
da abrupticidade do corte, e suavizar a transição resolve pouco. Em sinais
realistas o efeito é desprezível.

\paragraph{Efeito nas referências reais.} O que decide se as correções podem ser
ligadas por padrão é o desvio que produzem nos casos validados:

\begin{center}
\begin{tabular}{lcccccc|c}
\toprule
$R^A$ & T12/L1 & L1/L2 & L2/L3 & L3/L4 & L4/L5 & L5/S1 & desvio \\
\midrule
Atual (referência) & 0,98423 & 1,08503 & 1,22263 & 1,31024 & 1,28424 & 1,16975 & --- \\
Acolchoamento      & 0,98427 & 1,08508 & 1,22266 & 1,31034 & 1,28447 & 1,17030 & 0,047\% \\
Transição suave    & 0,98344 & 1,08249 & 1,21871 & 1,30432 & 1,27694 & 1,16272 & 0,602\% \\
As duas            & 0,98346 & 1,08248 & 1,21868 & 1,30429 & 1,27692 & 1,16266 & 0,607\% \\
\bottomrule
\end{tabular}
\end{center}

O acolchoamento desvia 0,047\%, abaixo do arredondamento dos valores publicados
--- todos os dígitos da referência continuam sendo reproduzidos. Corrige um
artefato de 30\% no pior caso a custo numérico nulo.

\textbf{Decisão:} apesar disso, o acolchoamento permanece \emph{desligado} por
padrão. A opção foi manter a cadeia idêntica à do software original, mesmo
conhecendo o artefato e seu custo. Quem precisar do resultado corrigido liga a
opção explicitamente e declara o fato.

A transição suave desvia 0,602\%, acima do arredondamento, e resolve pouco do
problema que pretendia atacar. Permanece como opção documentada, sem
recomendação de uso.

Os parâmetros são \code{core.ZERO\_PADDING} e \code{core.BAND\_TRANSITION}, e o
estudo é reproduzível por \code{tests/leakage\_study.py}.

\subsection{Interpolação da fase enrolada}
\label{sec:fasebug}

Este é o único defeito \emph{não} corrigido por padrão no núcleo do Anexo A, e a
razão é deliberada.

A função de transferência do modelo tem um salto de $2\pi$ na fase perto de
\unit{20}{Hz}. O software original interpola a fase enrolada, e a spline cúbica
atravessa o salto pelo caminho errado: nos cerca de 6\% de pontos em torno dele,
o valor complexo interpolado chega a errar 200\%.

Desenrolar a fase antes de interpolar corrige o problema, mas \emph{afasta o
resultado da referência}, porque o software original carrega o mesmo defeito:

\begin{center}
\begin{tabular}{lcccccc}
\toprule
$R^A$ & T12/L1 & L1/L2 & L2/L3 & L3/L4 & L4/L5 & L5/S1 \\
\midrule
Fase enrolada (padrão) & 0,98423 & 1,08503 & 1,22263 & 1,31024 & 1,28424 & 1,16975 \\
Fase desenrolada       & 0,98810 & 1,08905 & 1,22754 & 1,31568 & 1,29097 & 1,18076 \\
Diferença              & 0,39\% & 0,37\% & 0,40\% & 0,42\% & 0,52\% & 0,94\% \\
\bottomrule
\end{tabular}
\end{center}

O padrão é a fase enrolada, por fidelidade. A correção é acessível pela caixa
\emph{Unwrap transfer function phase before interpolating}, no painel
\emph{Parameters} da aba \emph{Annex A / Annex E} (\code{core.UNWRAP\_PHASE}).

\begin{quote}
\textbf{Ao publicar resultados, é indispensável declarar qual dos dois modos foi
usado.} O número que reproduz o software oficial é o da fase enrolada; o da fase
desenrolada é numericamente mais correto mas não corresponde ao que a ferramenta
de referência produz. A diferença de até 1\% é uma estimativa razoável do erro
numérico embutido no valor oficial.
\end{quote}

% =====================================================================
\section{Dados de campo do registrador do laboratório}
\label{sec:campo}

\subsection{Formato e características}

O registrador do laboratório grava arquivos \code{.m}, que são scripts MATLAB
contendo blocos de amostras precedidos de carimbo de tempo do relógio de tempo
real. O firmware confirma as características: MPU6050 na escala de $\pm 2\,g$,
portanto 16384 contagens por $g$; valores gravados como inteiros brutos, sem
nenhuma conversão ou calibração aplicada; 2048 amostras por bloco a
\unit{1024}{Hz} nominais, isto é, \unit{2}{s} de sinal; e um bloco a cada
\unit{5}{s} de relógio.

\textbf{O ciclo de trabalho é de 40\% e é intencional}: a cada \unit{5}{s},
\unit{2}{s} são gravados e \unit{3}{s} não. Um arquivo típico abrange 54 minutos
de relógio e contém 22 minutos de sinal. A taxa de amostragem efetiva de cada
bloco é medida pelo firmware e impressa no monitor serial, mas não é gravada no
arquivo; a taxa nominal é assumida.

\subsection{Consequência para os métodos da ISO 2631-5}

O caminho do Anexo A não se aplica a estes registros. A dose de sexta potência
depende de contar \emph{todos} os picos de compressão, e 60\% do tempo não foi
gravado --- a subestimativa resultante não é corrigível, porque os choques que
caíram nos intervalos simplesmente não existem para o cálculo.

Independentemente disso, os registros analisados apresentam picos de
aceleração de até \unit{21{,}3}{m/s^2}, mais que o dobro do critério de
\unit{9{,}81}{m/s^2}. Mesmo com gravação contínua, eles pertenceriam ao caminho
da Cláusula 5 com Anexo C. É fisicamente esperado: cadeira de rodas em piso
urbano perde contato ao descer guias e desníveis.

\paragraph{A Cláusula 5, ao contrário do Anexo A, não está descartada por
construção.} Ela usa só a direção $z$ no assento (Seção~2.3), então o
problema de sincronismo entre unidades da Seção~\ref{sec:conjuntoB} --- que
inviabiliza a soma multiponto da Fórmula (A.1) --- simplesmente não se aplica
a ela. O obstáculo que resta é só o ciclo de trabalho: cada bloco gravado tem
\unit{2}{s}, e o passa-alta de \unit{0{,}1}{Hz} da Seção 4 tem constante de
tempo da ordem de \unit{1{,}6}{s} --- levanta a dúvida de se o filtro chega a
se acomodar dentro de um bloco tão curto.

\code{tests/clause5\_block\_study.py} testa isso separando dois efeitos sobre
sinais sintéticos de resposta conhecida (choques com amplitude comparável aos
\unit{21{,}3}{m/s^2} reais, sobre vibração de fundo contínua): (i) a perda de
picos nos 60\% não gravados --- estatística, esperada, não corrigível pelo
processamento --- e (ii) o reinício do filtro a cada bloco independente, que
é a única parte passível de correção. Isolando (ii) sozinho --- comparando o
$D_z$ das mesmas amostras mantidas, processadas de uma vez versus em blocos
independentes --- o desvio ficou abaixo de 1\% em todas as durações e
comprimentos de bloco testados (2 a 20\,s), mesmo com o passa-alta
tecnicamente ainda em transiente durante o bloco inteiro pelo critério de
resposta ao degrau. Ou seja: o reinício do filtro a cada 2\,s, isoladamente,
não é o problema que a constante de tempo sugeria. O efeito (i) --- quanto os
picos amostrados representam o dia inteiro --- variou mais entre execuções
(de 82\% a 99\% do $D_z$ contínuo) e é uma questão de representatividade
estatística, não de implementação: só dados de campo reais, não um sinal
sintético, podem validá-la. Essa é a condição que falta para transformar isto
em funcionalidade do construtor de planilhas (Seção~13).

\subsection{Colagem dos blocos: um erro de amplitude de trinta vezes}
\label{sec:colagem}

A análise espectral anterior destes dados removeu os intervalos não gravados e
colou os blocos em um registro contínuo, aplicando janela e transformada sobre o
conjunto. A comparação com a análise por blocos, em que cada bloco de 2048
amostras é um segmento de Welch com janela própria, mostra uma diferença
sistemática:

\begin{center}
\begin{tabular}{lccc}
\toprule
\textbf{Banda} & \textbf{Colado} & \textbf{Por blocos} & \textbf{Razão} \\
\midrule
\unit{0{,}5}{} a \unit{5}{Hz}  & 0,0105 & 0,5499 & 0,02 \\
\unit{5}{} a \unit{20}{Hz}     & 0,0135 & 0,5290 & 0,03 \\
\unit{20}{} a \unit{50}{Hz}    & 0,0122 & 0,3798 & 0,03 \\
\unit{50}{} a \unit{100}{Hz}   & 0,0077 & 0,2564 & 0,03 \\
\bottomrule
\end{tabular}
\end{center}

\noindent
Valores eficazes de banda em \unit{}{m/s^2}, arquivo \code{1AA} das 10\,h, eixo
vertical.

A colagem subestima a energia em cerca de trinta vezes, de modo praticamente
uniforme entre bandas. A causa não é a descontinuidade entre blocos, como se
poderia supor, mas a janela: aplicada sobre um registro de 1,3 milhão de
amostras, ela atenua quase todo o sinal nas abas, e a correção de amplitude
calculada para o conjunto não recupera o que se perdeu em cada bloco.

Como o fator é aproximadamente constante entre bandas, \textbf{as formas
espectrais e a localização dos picos obtidas por colagem permanecem válidas}.
Os \textbf{valores absolutos de amplitude, não}: devem ser refeitos pela análise
por blocos. É por isso que o software \textbf{nunca cola os blocos}: cada bloco
de 2\,s é tratado como um segmento de Welch independente, na cadeia de cálculo e
no explorador de registros de campo (Seção~\ref{sec:fieldexplorer}).

\subsection{Calibração pela gravidade}

Com o sistema parado, o módulo do vetor medido deve ser a gravidade local,
\unit{9{,}784}{m/s^2} em Belo Horizonte. Isso dá uma referência absoluta e
gratuita para o ganho de cada unidade, e é a base da orientação automática de
eixo usada em todo o software (Seção~\ref{sec:eixos}). Os blocos iniciais do
registro \code{1AA} são comprovadamente estáticos --- desvio de
\unit{0{,}07}{m/s^2}, compatível com o piso de ruído do MPU6050 --- e fornecem
\unit{9{,}616}{m/s^2}, portanto um fator de correção de 1,0175.

Uma segunda unidade, \code{2BB}, forneceu \unit{10{,}74}{m/s^2}, um erro
aparente de escala de quase 10\%. A Seção~\ref{sec:conjuntoB} revê essa leitura:
o desvio é mais provavelmente \emph{offset}, não ganho, e a correção de ganho
por gravidade passou a vir desligada por padrão.

\subsection{Avaliação pela NHO-09}

A NHO-09 da Fundacentro combina os três eixos por soma vetorial, com fatores
1,4 em $x$ e $y$ e 1,0 em $z$, sobre acelerações ponderadas por $W_d$ e $W_k$.
Difere da diretiva europeia, que usa o eixo de maior valor. Os critérios são
$a_{ren}$ de \unit{0{,}5}{m/s^2} e $VDVR$ de \unit{9{,}1}{m/s^{1{,}75}} para
o nível de ação, e \unit{1{,}1}{m/s^2} e \unit{21}{m/s^{1{,}75}} para o limite
de exposição. Basta um dos dois critérios ser ultrapassado.

\paragraph{A média aritmética prescrita subestima em 70\%.} Para cada
componente de exposição, a norma prescreve a média \emph{aritmética} dos
valores $a_{mr}$ das repetições. Isso é adequado para repetições da mesma
operação, com amplitudes semelhantes. Aplicada aos blocos de um dia de rotina,
com longos trechos parados e poucos blocos intensos, a média aritmética deu
\unit{0{,}110}{m/s^2} contra \unit{0{,}360}{m/s^2} da média pela energia ---
69,6\% abaixo. Aplicada literalmente, a fórmula produziria um falso veredito de
segurança. A implementação trata o dia como uma única componente amostrada
sistematicamente, agregando pela energia, e reporta o valor aritmético apenas
como verificação.

\paragraph{Invariância à ocupação.} Parte do dia medido corresponde,
provavelmente, à cadeira desocupada: por quatro horas seguidas a orientação do
assento ficou constante com desvio de centésimos de grau, enquanto nos períodos
ativos ela oscila de 1 a 2 graus. O resultado da NHO-09, porém, é idêntico
incluindo ou excluindo esses trechos --- $a_{ren} = \unit{0{,}369}{m/s^2}$ e
$VDVR = \unit{12{,}62}{m/s^{1{,}75}}$ nos dois casos. Os blocos parados
contribuem quase nada de energia, e excluí-los reduz o tempo de exposição na
mesma proporção, de modo que os efeitos se cancelam. A condição é que os blocos
usados e o tempo declarado venham do mesmo período. A detecção de ocupação é
hoje uma informação exibida na interface, não um filtro aplicado
(Seção~\ref{sec:fieldlogger}).

\paragraph{Critérios discordantes.} Para o dia analisado, o $a_{ren}$ fica
abaixo do nível de ação e o $VDVR$ acima dele. É coerente com o conteúdo de
choque do sinal e com a literatura, que aponta o VDV como o critério mais
restritivo em exposições com choques.

\paragraph{Extrapolação para jornadas longas.} A normalização
$a_{ren} = a_{re}\sqrt{T/T_0}$ já é o mecanismo que dá conta de exposições
maiores que 8 horas, e o VDV acumula ao longo do tempo real. O que não está
estabelecido é a hipótese de equivalência por energia fora da faixa de 4 a 8
horas: o Anexo B da ISO 2631-1 apresenta uma segunda relação, com raiz quarta
do tempo, e as duas divergem para durações longas. Extrapolando o dia medido
para 16 horas, a relação de raiz quadrada dá \unit{0{,}509}{m/s^2}, acima do
nível de ação, e a de raiz quarta dá \unit{0{,}428}{m/s^2}, abaixo. O software
reporta as duas, e a extrapolação exige escolha explícita do usuário na
interface (caixa \emph{Extrapolate}, Seção~\ref{sec:fieldlogger}).

\subsection{Avaliação pela Diretiva 2002/44/CE}

A diretiva distingue explicitamente dois casos. Para vibração mão-braço,
prescreve a raiz da soma dos quadrados dos três eixos; para corpo inteiro,
prescreve o \emph{maior} dos três valores ponderados $1{,}4\,a_{wx}$,
$1{,}4\,a_{wy}$ e $a_{wz}$, tanto para o $A(8)$ quanto para o VDV. Os valores são
\unit{0{,}5}{m/s^2} ou \unit{9{,}1}{m/s^{1{,}75}} para ação e
\unit{1{,}15}{m/s^2} ou \unit{21}{m/s^{1{,}75}} para limite. O método do
laboratório soma os eixos também em corpo inteiro, citando o guia de boas
práticas da Comissão Europeia de 2007, o que é uma opção mais conservadora mas
diverge do texto legal. O software oferece os dois modos, através do seletor
\emph{Criterion}. Registre-se ainda que o artigo 4(4)(g) da diretiva lista o
prolongamento da exposição além do horário normal de trabalho como fator a
considerar na avaliação --- respaldo direto para a discussão sobre jornadas
longas.

\paragraph{Soma vetorial e NHO-09 coincidem.} A soma vetorial dos $A(8)$ parciais,
$\sqrt{\sum_i (k_i\,a_{wi}\sqrt{t/8})^2} = \sqrt{t/8}\,\sqrt{\sum_i (k_i\,a_{wi})^2}$,
é algebricamente idêntica ao $a_{ren}$ da NHO-09. Nos dados do assento, as duas
dão \unit{0{,}369}{m/s^2}.

\paragraph{Validação contra a calculadora do HSE.} A lógica de agregação do $A(8)$
foi conferida contra a calculadora oficial de vibração de corpo inteiro do HSE
britânico, versão 1.06. Para as quatro operações de exemplo da planilha, os
pontos de exposição por eixo coincidem exatamente (105,5 / 149 / 35,4), e o
valor diário é o do maior eixo, \unit{0{,}61}{m/s^2}, como prescreve a diretiva.
A soma vetorial dos mesmos dados daria \unit{0{,}85}{m/s^2}.

\paragraph{O guia de boas práticas prescreve o maior eixo.} O método do
laboratório justifica a soma dos três eixos citando o guia de boas práticas da
Comissão Europeia. O próprio guia, no Anexo E.4, prescreve para o VDV: combinar
os VDV parciais de cada \emph{tarefa} por eixo, e tomar como valor diário o
\emph{maior} dos três eixos. A soma é entre tarefas, não entre eixos.

\subsection{Auditoria do software de análise do laboratório}
\label{sec:auditoria}

O código-fonte do software em MATLAB usado nos trabalhos do laboratório foi
inspecionado: o conversor \code{ConverteDados.m}, a ponderação
\code{Gera\_Isofil.m} e o aplicativo \code{WBV\_HSE\_v16\_Usuario}. O pipeline foi
reproduzido em Python sobre os mesmos arquivos, e permanece disponível na
interface como o painel \emph{Lab method} (Seção~\ref{sec:posicao}), para
comparação lado a lado com este software.

\paragraph{Retratação de duas hipóteses anteriores.} Uma versão anterior desta
documentação sugeria que o VDV publicado retinha a gravidade. \textbf{Isso está
errado}: o conversor subtrai a média de cada bloco em cada eixo antes de
qualquer cálculo. Sugeria também que o $A(8)$ usava apenas o tempo gravado,
subestimando em 37\%. \textbf{Também está errado para o código}: o tempo usado é
a diferença entre o primeiro e o último carimbo de cada arquivo, isto é, o
período coberto. O texto do artigo, ao descrever o tempo como ``a soma do tempo
com dado válido'', descreve o cálculo de forma imprecisa; o código faz o
correto.

\paragraph{O VDV omite o passo de tempo.} A função \code{CalculaMeasuredVDV}
calcula
\[
\mathrm{mVDV} = \Big(\sum a_w^4\Big)^{1/4},
\qquad\text{em vez de}\qquad
\mathrm{VDV} = \Big(\sum a_w^4\,\Delta t\Big)^{1/4}.
\]
Com $\Delta t = 1/1024$\,s, o VDV resulta $1024^{1/4} = 5{,}657$ vezes maior que
o valor correto. A reprodução fiel do pipeline sobre os dados do assento
confirma a razão com precisão de quatro algarismos. Este defeito afeta todo VDV
produzido pelo software do laboratório, independentemente do conjunto de dados.
Este software calcula o VDV sempre pela forma correta, com o passo de tempo; a
caixa \emph{Reproduce the app VDV defects} existe apenas para conferir números
contra a tela do aplicativo original quando necessário.

\paragraph{O VDV parcial subestima a extrapolação.} A função
\code{CalculaPartialVDV} usa exposição fixa de uma hora e, como tempo medido, o
período coberto do arquivo. Mas o VDV medido é somado apenas sobre as amostras
gravadas, que são 40\% do período. O fator correto seria
$(T_{\mathrm{exp}}/T_{\mathrm{gravado}})^{1/4}$; o usado corresponde a
$(T_{\mathrm{exp}}/T_{\mathrm{coberto}})^{1/4}$, 20\% menor. Os dois defeitos se
compensam parcialmente: o VDV produzido fica cerca de 4,5 vezes acima do
correto.

\paragraph{Consequência para o artigo.} O VDV de \unit{146{,}3}{m/s^{1{,}75}}
reportado para o assento, corrigido apenas pelo passo de tempo, cai para cerca
de \unit{26}{m/s^{1{,}75}}. Se os dados forem de outra coleta --- a estrutura de
pastas do código sugere um conjunto ``B'' distinto do analisado aqui ---, esse
valor corrigido ainda estaria acima do limite de exposição. A conclusão
qualitativa pode sobreviver; os números, não.

\paragraph{O que está correto.} O valor eficaz ponderado é agregado pela média
quadrática dos subintervalos, que é a forma correta, e não pela aritmética. O
$A(8)$ usa o período coberto. A ponderação $W_k$ e $W_d$ reproduz a norma. As
falhas de leitura I$^2$C, que o firmware grava como $-1$, são substituídas pela
média do bloco --- tratamento que este software não fazia e passou a fazer,
seguindo o do laboratório.

\paragraph{Uma correção em sentido contrário.} A ponderação $W_b$ do código do
laboratório usa os parâmetros da ISO 8041, adotados pela ISO 2631-4: $Q_4 = 0{,}55$,
$Q_5 = 0{,}9$, $Q_6 = 0{,}95$ e ganho $K = 1{,}024$. Este software usava
$Q_4 = 0{,}63$ e $Q_5 = Q_6 = 0{,}8$, valores sem correspondência normativa, e foi
corrigido. Registre-se que o código do laboratório define o ganho $K$ mas não o
aplica.

\paragraph{Filtragem por bloco.} Os dois softwares ponderam cada bloco de
\unit{2}{s} isoladamente. O passa-alta de \unit{0{,}4}{Hz} da ponderação tem
período de \unit{2{,}5}{s}, maior que o bloco: o filtro não atinge regime
permanente dentro de um bloco. É limitação do dado, não de nenhum dos códigos, e
afeta a precisão abaixo de cerca de \unit{1}{Hz}.

\subsection{Conjunto B: campanhas de 2024 e 2025}
\label{sec:conjuntoB}

O conjunto B reúne as campanhas de outubro de 2024 a março de 2025, com arquivos
brutos das três posições. Excluídos os arquivos de teste e os datados por relógio
reiniciado --- 07/04/2024 às 0\,h é o valor que o próprio firmware grava após
reinício, e datas em 2020, 2046 e 2099 indicam relógio sem bateria ---, restam
134 arquivos.

\paragraph{Rótulos.} O arquivo rotulado como encosto no dia 17/10/2024 coincide
amostra por amostra com o registro \code{2BB} do conjunto A, com correlação
unitária. Os rótulos originais estão corretos. A hipótese de troca, levantada
anteriormente a partir da orientação da gravidade, está refutada: a orientação
do sensor dentro de cada base não é uniforme.

\paragraph{Orientação.} O eixo que carrega a gravidade muda entre campanhas para
uma mesma posição --- o encosto aparece com a gravidade em $z$, em $-z$ e em
$-x$ conforme a data --- e chega a mudar dentro de um mesmo dia: em 17/10/2024 o
sensor do encosto passa de $-z$ para $-y$ às 18\,h, indicando que a fixação se
soltou ou foi refeita. Nenhuma regra fixa de troca de eixos é aplicável; a
orientação precisa ser determinada arquivo a arquivo pela gravidade. É esse
achado que motivou a detecção automática de eixo por gravidade, hoje central em
todo o software (Seção~\ref{sec:eixos}).

\paragraph{Ganho ou offset: revisão de uma recomendação anterior.} Esta
documentação recomendava corrigir o ganho de cada unidade pela razão entre a
gravidade medida e a local. \textbf{A recomendação é revista.} Com uma única
orientação, um desvio do módulo medido pode ser erro de ganho ou erro de zero
(offset), e os dois são indistinguíveis. Três evidências apontam para offset:
uma mesma unidade dá desvios diferentes conforme o eixo que carrega a gravidade
(+8\% com a gravidade em $z$, +3,6\% em $y$); o desvio de +8\% corresponde a
\unit{0{,}78}{m/s^2}, exatamente o limite de offset do eixo $z$ na folha de
dados do MPU6050, 80\,mg; e desvios de até 5\% em ganho estariam fora da
tolerância de sensibilidade do componente, enquanto como offset estão dentro
dela. Offset não afeta a medição de vibração, porque a componente estática é
removida; corrigi-lo como se fosse ganho introduziria erro. A correção passa a
vir desligada por padrão. Retira-se, em particular, a afirmação anterior de que
o desvio de 9\% do sensor do encosto implicaria 77\% de erro na dose.

A separação correta exige calibração em seis posições, com cada eixo apontado
para cima e para baixo: o offset é a média das duas leituras e o ganho, a
semidiferença.

\paragraph{Sincronismo entre sensores.} Os blocos das três posições são
pareáveis pelo carimbo de tempo, porque todas as unidades disparam nos múltiplos
de \unit{5}{s} do próprio relógio. O alinhamento dentro do bloco, porém, não
existe. Na sessão de 25/03/2025, a correlação cruzada entre assento e encosto
nos blocos com movimento dá atrasos espalhados de $-1{,}2$ a $+1{,}0$\,s, sem
tendência consistente. Cada relógio de tempo real vira o segundo num instante
diferente, e o atraso relativo varia de bloco a bloco.

As consequências são diretas. A avaliação por posição, pela NHO-09 ou pela
diretiva, não é afetada. A soma entre pontos de contato da Fórmula (A.1) é
inviável, o que é uma segunda razão, independente do ciclo de trabalho, para o
Anexo A não se aplicar a estes dados. E a transmissibilidade entre posições só
pode ser estimada em módulo, pela razão entre autoespectros; fase e coerência
exigem relógio comum. O próximo registrador precisa de sincronismo em hardware.

\paragraph{Resultados.} Com o tempo de exposição igual ao período coberto e sem
correção de ganho:

\begin{center}
\small
\begin{tabular}{llrrrrr}
\toprule
Sessão & Posição & Período & $a_{ren}$ & $VDVR$ & $A(8)$ máx & VDV máx \\
\midrule
17/10/2024 & encosto & 8,39\,h & 0,441 & \textbf{17,96} & 0,410 & \textbf{17,90} \\
25/10/2024 & assento & 9,58\,h & 0,349 & \textbf{12,84} & 0,223 & \textbf{11,47} \\
25/10/2024 & encosto & 8,85\,h & 0,300 & \textbf{9,94}  & 0,206 & 8,98 \\
04/11/2024 & encosto & 7,38\,h & 0,049 & 6,35  & 0,033 & 6,30 \\
20/01/2025 & assento & 8,38\,h & 0,330 & 8,79  & 0,218 & 8,28 \\
20/01/2025 & encosto & 8,54\,h & 0,416 & \textbf{10,54} & 0,255 & \textbf{9,30} \\
25/03/2025 & assento & 7,23\,h & 0,271 & \textbf{10,26} & 0,225 & \textbf{10,12} \\
25/03/2025 & encosto & 7,23\,h & 0,366 & \textbf{12,43} & 0,274 & \textbf{11,91} \\
\bottomrule
\end{tabular}
\end{center}

\noindent
Em negrito, acima do nível de ação. O padrão é consistente em todas as sessões:
o valor eficaz fica abaixo do nível de ação, e o VDV fica acima dele na maioria
das sessões ativas, sempre abaixo do limite de exposição. É o perfil de uma
exposição dominada por choques. O resultado de 17/10 no encosto deve ser lido
com reserva, porque o sensor se soltou naquele dia e o deslocamento da fixação
produz acelerações que não são vibração transmitida ao usuário.

\subsection{Reconstrução do resultado do encosto publicado}
\label{sec:encosto}

O trabalho de análise destes dados reporta, para o encosto, $A(8)$ de
\unit{0{,}685}{m/s^2}, acima do nível de ação. A configuração usada foi
recuperada: arquivos de 17/10/2024, das 10\,h às 18\,h, e um arquivo de
18/10/2024 às 11\,h; ponderação $W_c$ no eixo $x$ com $k = 0{,}8$, nenhuma
ponderação em $y$ e $z$ com $k = 1$; e a opção de troca dos eixos $x$ e $z$
ligada. Reproduzindo essa configuração, os trinta valores de $a_w$ por arquivo
e todas as parcelas de $A(8)$ coincidem com a saída do software do laboratório
na terceira casa, e a soma vetorial dos totais dá \unit{0{,}685}{m/s^2}.

A decomposição mostra de onde vem o número:

\begin{center}
\small
\begin{tabular}{lcccc}
\toprule
Passo & $x$ & $y$ & $z$ & Combinado \\
\midrule
1. Configuração publicada                        & 0,299 & 0,428 & 0,442 & 0,685 \\
2. Sem o arquivo de 18/10, de outro dia           & 0,269 & 0,391 & 0,412 & 0,628 \\
3. Com a orientação correta, sem troca de eixos   & 0,147 & 0,391 & 0,606 & 0,736 \\
4. Com $y$ e $z$ ponderados                       & 0,147 & 0,037 & 0,057 & 0,162 \\
\bottomrule
\end{tabular}
\end{center}

Cinco problemas se acumulam:

\begin{enumerate}
\item \textbf{Os eixos $y$ e $z$ entraram sem ponderação.} A opção ``nenhuma''
      do software não aplica sequer a limitação de banda: o valor eficaz inclui
      tudo de 0 a \unit{512}{Hz}, ruído do sensor inclusive. É a principal causa
      do resultado --- ponderados, os dois eixos caem de 0,391 e 0,606 para
      0,037 e 0,057.
\item \textbf{A troca de eixos estava errada para esta sessão.} Neste dia a
      gravidade está na coluna 3 do sensor, que portanto é o eixo vertical $z$.
      A troca levou esse eixo para a posição $x$ e aplicou a ele a ponderação
      $W_c$, própria da direção ântero-posterior.
\item \textbf{Um arquivo de outro dia entrou na exposição diária.} O $A(8)$ é
      uma grandeza diária; o registro de 18/10 não pertence à jornada de 17/10.
\item \textbf{O encosto não integra a avaliação de saúde.} A nota da subcláusula
      7.2.3 da ISO 2631-1 recomenda medir o eixo $x$ do encosto com $W_c$ e
      $k = 0{,}8$, mas registra que, pela falta de evidência sobre o efeito
      desse movimento na saúde, ele não entra na avaliação de severidade do
      Anexo B. Os eixos $y$ e $z$ do encosto não têm ponderação de saúde. A
      comparação com o nível de ação da diretiva não tem, portanto, respaldo
      normativo. O único valor de saúde calculável é o $A(8)$ em $x$, de
      \unit{0{,}147}{m/s^2}, e ele é informativo.
\item \textbf{O registro das 17\,h é suspeito.} É o de maior valor do dia e
      antecede a mudança de orientação do sensor às 18\,h. Excluindo-o, o
      $A(8)$ em $x$ cai para \unit{0{,}097}{m/s^2}.
\end{enumerate}

A conclusão de que o encosto ultrapassou o nível de ação não se sustenta.

\paragraph{Salvaguarda no software.} Para impedir que o mesmo erro se repita, a
avaliação passou a depender da posição do sensor, declarada pelo usuário
(Seção~\ref{sec:posicao}). No assento, avaliação de saúde. No encosto, avaliação
de conforto pelo valor total de vibração $a_v$, com $W_c$, $W_d$ e $W_d$ e
fatores 0,8, 0,5 e 0,4, e o eixo $x$ com $W_c$ e $k = 0{,}8$ exibido apenas como
valor informativo, sem comparação com nível de ação. Nos pés, avaliação de
conforto com $W_k$ e fatores 0,25, 0,25 e 0,4. Aplicado aos mesmos arquivos do
artigo, o tratamento de assento reproduz o veredito publicado; o de encosto dá
$a_v = \unit{0{,}186}{m/s^2}$, faixa ``não desconfortável'' do Anexo C.

Registram-se ainda duas observações sobre o software do laboratório. Os totais
por eixo exibidos na tela analisada, 0,220 / 0,428 / 0,422, não correspondem à
soma quadrática das parcelas exibidas na mesma tela, 0,299 / 0,428 / 0,442; o
resultado publicado é coerente com as parcelas, o que sugere que o painel de
totais refletia outro estado da opção de troca de eixos, que é mantida arquivo
a arquivo. E o VDV diário é combinado entre eixos pela raiz da soma dos
quadrados, operador que não corresponde nem ao maior eixo prescrito pela
diretiva nem à soma de quartas potências própria do VDV.

\subsection{Eixos}
\label{sec:eixos}

Nos dados do assento a gravidade aparece na primeira coluna do arquivo, e nos
dados do encosto na terceira --- o inverso da descrição inicial da montagem. O
processamento segue os dados: a coluna que contém a gravidade é tratada como
eixo vertical, arquivo a arquivo, sem regra fixa por posição. Esta é hoje a
detecção padrão em todo o software; a antiga opção manual de ``troca x-z''
sobrevive apenas como caixa de auditoria \emph{Force fixed x-z swap} (modo
\emph{audit}), desligada por padrão, para reproduzir o comportamento antigo
quando necessário.

% =====================================================================
\section{Do registro forense ao painel produtizado}
\label{sec:fieldlogger}

As Seções~\ref{sec:campo} a~\ref{sec:eixos} descrevem um trabalho de forense de
dados: reconstrução manual de um resultado publicado, cruzamento de arquivo por
arquivo, decisões tomadas caso a caso. Esse trabalho foi incorporado ao
software como dois recursos permanentes, descritos aqui antes da arquitetura
geral (Seção~\ref{sec:arquitetura}) porque decorrem diretamente das descobertas
acima.

\subsection{Painel Field logger / NHO-09}

Dentro do explorador de sinais, a sub-aba \emph{Field logger / NHO-09} carrega
um dia inteiro de arquivos \code{.m} de uma unidade, aplica a calibração pela
gravidade e a detecção automática de eixo, e produz lado a lado: a avaliação
NHO-09 (Seção~\ref{sec:campo}), a avaliação pela Diretiva 2002/44/CE pelo maior
eixo e pela soma vetorial, e o painel \emph{Lab method}, que reproduz o
aplicativo \code{WBV\_HSE\_v16} com a configuração escolhida pelo usuário ---
ponderação por eixo, fatores $k$, troca x-z, base de tempo do valor eficaz.
Um quadro no formato do painel \emph{Total A(8) Exposure} do aplicativo original
mostra o $A(8)$ de cada eixo e o diário, pelo maior eixo e pela soma vetorial,
em qualquer posição do sensor, com as duas linhas --- este software e o método
do laboratório --- lado a lado.

\label{sec:posicao}
O seletor \emph{Sensor position} aplica o tratamento normativo correto por
ponto de contato (assento: saúde; encosto: conforto, com o eixo $x$
informativo; pés: conforto), a mesma regra descrita na
Seção~\ref{sec:encosto}. Três gráficos comparam este software com o método do
laboratório: um \emph{dashboard} com cada critério em sua própria escala e
faixas verde/âmbar/vermelha; uma comparação \emph{side by side} eixo a eixo com
diferença percentual, pareada pelo eixo físico do sensor e não pela letra do
rótulo; e a contribuição de cada arquivo à energia do $A(8)$ diário, que mostra
onde está a exposição do dia.

\subsection{Explorador de registros de campo}
\label{sec:fieldexplorer}

Uma segunda sub-aba, \emph{Field logger records (.m)}, inspeciona os mesmos
arquivos \code{.m} independentemente do fluxo NHO-09/Diretiva --- não calcula
dose nem $A(8)$, apenas caracteriza o sinal. Três visualizações: linha do tempo
de RMS por bloco, espectro por bloco (Welch, com a opção de janela CONEM usada
no laboratório) e forma de onda de um único bloco, selecionável por índice.
Existe porque a triagem de um registro --- verificar se um sensor se soltou,
como no dia 17/10/2024 (Seção~\ref{sec:conjuntoB}), ou se a orientação mudou no
meio do dia --- é uma necessidade recorrente e independente de qualquer cálculo
normativo.

A tabela de estatísticas por arquivo usa \code{fielddata.field\_statistics}:
número de blocos, tempo gravado, extensão, ciclo de trabalho, RMS por eixo e
total, pico, fator de crista, VDV, curtose, frequência dominante por eixo
(via espectro de Welch por bloco), eixo de gravidade e mapa de eixos.

% =====================================================================
\section{Análise de sensibilidade postural e antropométrica}
\label{sec:sensibilidade}

Este recurso nasceu diretamente da recomendação metodológica da
Seção~\ref{sec:dados}: como nenhuma das sete posturas tabeladas descreve um
usuário de cadeira de rodas, e como o índice antropométrico de 30 posições não
está univocamente decifrado (Seção~\ref{sec:indice}), a dispersão do resultado
entre todas as combinações é, em si, o resultado defensável.

A função \code{core.sensitivity\_grid} varre as $7$ posturas disponíveis
(Seção~\ref{sec:dados}) contra os $30$ modelos antropométricos, para um mesmo
conjunto de excitações medidas, reaproveitando o cache de resultado leve
descrito na Seção~9.3 para manter as 210 combinações viáveis em memória. Cada
combinação passa pela cadeia completa do Anexo A/Anexo E (Seção~\ref{sec:cadeiaA}),
e o resultado é um par de arranjos $S_d^A$ e $R^A$ por nível discal, postura e
antropometria.

A interface apresenta um gráfico de caixas por nível discal, cobrindo as 210
combinações, com a combinação efetivamente medida (postura e antropometria
declaradas na planilha) destacada sobre a distribuição. Isso torna visível, para
o mesmo sinal medido, o quanto o resultado normativo depende de premissas
biomecânicas que não foram de fato verificadas para o sujeito --- particularmente
relevante para populações fora do padrão de operador adulto de máquina pesada,
como os usuários de cadeira de rodas que motivam este trabalho, e, por extensão,
qualquer futura aplicação a crianças (Seção~\ref{sec:criancas}).

Esse recurso, porém, exige um sinal no formato de excitação do Anexo~A --- e
dados de campo do registrador não podem entrar nesse formato, pelos dois
motivos já documentados na Seção~\ref{sec:campo} (ciclo de trabalho de 40\% e
ausência de sincronismo entre unidades). Isso deixava justamente a população
que motivou o recurso --- usuários de cadeira de rodas medidos em campo --- sem
como usá-lo. Esta versão adiciona um equivalente para a Cláusula~5/Anexo~C, na
própria aba \emph{Clause 5 / Annex C} (gráfico \emph{Anthropometric
sensitivity}, disponível tanto para registro ASCII quanto para importação de
campo): como a Cláusula~5 não tem parâmetro de postura --- a função de
transferência assento-coluna da Fórmula~1 é fixa ---, a dispersão varrida é só
antropométrica, sobre sexo, fração de massa acima da pelve e área de placa
terminal, cada um a $\pm 20\%$ do valor em uso. A Fórmula~(5) ($D_{zd}$) não
depende de nenhum dos três, então a varredura reaproveita o $D_{zd}$ já
calculado e recalcula só $m_z$, $S_d$, $R$ e $\Pi$ (Fórmulas~C.2, C.3, C.5)
para cada combinação --- 50 valores de $R$ em vez de reprocessar o sinal 50
vezes, o que a torna instantânea mesmo sobre um dia inteiro de blocos de
campo. O gráfico de caixas original do Anexo~A/E permanece exatamente onde
está; este é um recurso irmão, não uma substituição, porque as duas cadeias
varrem dimensões de incerteza diferentes (postura+modelo vs. só antropometria).

% =====================================================================
\section{Arquitetura do software}
\label{sec:arquitetura}

\begin{center}
\small
\begin{tabular}{ll}
\toprule
\textbf{Módulo} & \textbf{Responsabilidade} \\
\midrule
\code{core.py}            & Anexo A/E, ponderações ISO 2631-1, métodos espectrais, análise de sensibilidade \\
\code{annexc.py}          & Cláusula 5 e Anexo C, independente de \code{core} \\
\code{workbook.py}        & leitura e escrita das planilhas de exposição \\
\code{analysis.py}        & orquestração das análises orientadora e individual \\
\code{qtworker.py}        & \code{Worker} e \code{run\_in\_thread}: infraestrutura de execução em segundo plano \\
\code{gui.py}             & janela principal, orquestra as abas do avaliador \\
\code{clause5\_tab.py}    & aba da Cláusula 5 \\
\code{explorer.py}        & explorador de sinais ASCII (independente da cadeia normativa) \\
\code{fielddata.py}       & leitura de \code{.m}, calibração pela gravidade, orientação, estatísticas \\
\code{field\_panel.py}    & painel Field logger / NHO-09, em segundo plano \\
\code{field\_explorer.py} & explorador de registros \code{.m} (Seção~\ref{sec:fieldexplorer}) \\
\code{nho09.py}           & avaliação NHO-09, cache de blocos ponderados \\
\code{directive.py}       & avaliação pela Diretiva 2002/44/CE \\
\code{labmethod.py}       & reprodução auditada do \code{WBV\_HSE\_v16} do laboratório \\
\code{positions.py}       & tratamento normativo por posição de sensor (assento/encosto/pés) \\
\code{summary.py}         & agregação dos critérios (NHO-09, laboratório, por arquivo) \\
\code{weighting\_panel.py}& painel de ponderação por eixo \\
\code{editor.py}          & construtor de planilhas de histórico de exposição \\
\code{theme.py}           & paleta e folha de estilo, compartilhada pelos dois executáveis \\
\code{build\_data.py}     & conversão do pacote de dados original \\
\bottomrule
\end{tabular}
\end{center}

O núcleo de cálculo pode ser usado sem a interface, em análises em lote ou
estudos paramétricos, por importação direta dos módulos.

\subsection{Execução em segundo plano}
\label{sec:threading}

Carregar dezenas de arquivos \code{.m} ou recalcular a NHO-09 ao trocar de
postura envolve, tipicamente, dezenas de segundos de processamento --- a leitura
de um dia completo de arquivos chega a 41\,s, e a NHO-09 normativa, a 66\,s,
para um conjunto de 81 arquivos e cerca de 35 mil blocos. Feito na thread da
interface, isso travava a janela pelo tempo do cálculo.

O padrão adotado é \code{Worker}, um \code{QObject} movido para uma
\code{QThread} por \code{qtworker.run\_in\_thread}, com sinais conectados
sempre a métodos vinculados (nunca a \emph{lambdas}, que em certas condições o
PySide6 falha em associar corretamente à thread de origem --- um defeito de
concorrência real, encontrado e corrigido nesta versão, descrito adiante). Cada
painel expõe uma barra de progresso e um rótulo de status, desabilita os
controles que dependem do resultado em andamento, e reconecta tudo ao final,
com sucesso ou falha.

\paragraph{Defeito de concorrência corrigido.} Uma primeira versão do painel de
campo conectava os sinais de progresso e conclusão da \code{QThread} a funções
\emph{lambda}. Isso produzia, de forma intermitente, travamento completo do
processo ao carregar 10 arquivos e trocar de postura em sequência rápida.
Instrumentação com \code{faulthandler.dump\_traceback\_later} revelou que o
código do \emph{callback} às vezes executava na thread de \emph{worker}, não na
thread da interface --- o PySide6 depende do tipo de conexão para decidir a
afinidade de thread, e essa decisão pode falhar quando o alvo é uma
\emph{lambda} em vez de um método vinculado de um \code{QObject} com afinidade
conhecida. A correção foi conectar todos os sinais diretamente a métodos
vinculados, guardando o estado necessário (por exemplo, qual função de retomada
chamar) como atributo de instância em vez de capturá-lo por \emph{closure}. O
cenário de falha original --- carregar 10 arquivos, trocar de posição três vezes
em sequência --- foi repetido múltiplas vezes após a correção, sempre com
sucesso.

Duas proteções adicionais acompanham o padrão: um sinalizador de reentrância
impede que uma segunda operação seja disparada enquanto a primeira ainda roda
(o que arriscaria substituir a referência à \code{QThread} em andamento e deixar
o objeto órfão em execução); e todo caminho de finalização, com sucesso ou
falha, chama \code{thread.quit()} seguido de \code{thread.wait()}, evitando o
vazamento de threads que a primeira versão do painel apresentava.

\subsection{Formato de entrada}

O software lê as mesmas planilhas do software original de referência, com três
abas: exposições (identificador, grupo de postura, duração da medição e os
arquivos dos quatro pontos de contato), padrões diários (combinações de
exposições com durações) e anos (massa corporal, estatura e os padrões
praticados com o número de dias). O aplicativo auxiliar \emph{Avaliador de
Vibração Humana --- Planilhas} (\code{editor.py}) constrói essas planilhas sem
edição manual, detectando automaticamente a duração de cada arquivo de
exposição e mostrando, para cada combinação de massa e altura, em qual dos 30
modelos antropométricos ela cai.

\subsection{Abas do avaliador}

\begin{description}
\item[Annex A / Annex E] parâmetros da análise: planilha de entrada, pasta de
      acelerações, opções de fase enrolada/desenrolada
      (Seção~\ref{sec:fasebug}), janela e preenchimento com zeros
      (Seção~\ref{sec:leakage}), e logotipo do GRAVI-SH. Renomeada nesta
      versão (era \emph{Setup}) para deixar explícito que ela configura
      especificamente o caminho do Anexo A/Anexo E, e não o programa como um
      todo --- espelhando o nome já usado pela aba \emph{Clause 5 / Annex C}.
\item[Results] tabela de $S_d^A$ e $R^A$ por nível discal, aviso de validade
      dos \unit{9{,}81}{m/s^2} acima da tabela (verde ou vermelho).
\item[Frequency domain] espectro das excitações, com seleção de ponto de
      contato e eixo, sobreposição do espectro não processado, o painel de
      ponderação em frequência, e --- nesta versão --- fase e coerência
      (Seção~\ref{sec:fasecoerencia}).
\item[Time domain] força de compressão por nível discal, com histograma dos
      picos contados pela Equação~\eqref{eq:a2}.
\item[Signal processing] tabela completa dos parâmetros de processamento
      aplicados a cada excitação: amostras, taxa, comprimento de transformada,
      janela, resolução em frequência, Nyquist, cortes de banda e bins
      restantes.
\item[Clause 5 / Annex C] segundo caminho da norma (Seção~\ref{sec:contato}),
      independente das demais abas. Além do registro ASCII contínuo original,
      esta versão aceita um dia inteiro de arquivos \code{.m} do registrador de
      campo (Seção~\ref{sec:licoes}): cada bloco de \unit{2}{s} é condicionado e
      tem a perda de contato detectada de forma independente, sem colagem, e o
      tempo medido conta só os \unit{2}{s} de fato gravados por bloco, não os
      \unit{5}{s} do ciclo. Cinco gráficos: sinal condicionado e resposta da
      coluna, histograma de picos, função de transferência, espectro em
      frequência do $a_z$ bruto (informativo --- a Cláusula 5 não pondera por
      frequência), as faixas de referência de $R$/$\Pi$ contra a Tabela~C.2, e
      a sensibilidade antropométrica (Seção~\ref{sec:sensibilidade}) --- os
      três adicionados nesta versão, os dois primeiros motivados pela revisão
      de literatura da Seção~\ref{sec:licoes}.
\item[Signal explorer] três sub-abas: registros ASCII (\code{explorer.py}),
      Field logger / NHO-09 (Seção~\ref{sec:fieldlogger}) e Field logger
      records .m (Seção~\ref{sec:fieldexplorer}).
\item[Sensitivity analysis] análise de sensibilidade postural e antropométrica
      (Seção~\ref{sec:sensibilidade}).
\end{description}

As abas \emph{Results}, \emph{Frequency domain}, \emph{Time domain} e
\emph{Signal processing} só aparecem depois da primeira análise; ao abrir, o
programa mostra apenas \emph{Annex A / Annex E}, \emph{Clause 5 / Annex C} e
\emph{Signal explorer}.

\subsection{Recursos de visualização e interface}

O explorador de sinais ASCII analisa arquivos diretamente, sem planilha,
sobrepondo várias séries com caixas para ligar e desligar cada uma, em três
visualizações: série temporal, espectro (cinco métodos e cinco janelas,
Seção~\ref{sec:espectral}) e RMS acumulado, útil para verificar
estacionariedade. A tabela lateral fornece, por arquivo: número de amostras,
duração, taxa de amostragem, RMS por eixo e soma vetorial, pico, fator de
crista, VDV, curtose, frequência dominante e a classe de intensidade que a
norma atribuiria.

A curtose merece nota: vale 3 para ruído gaussiano, e valores muito acima
indicam conteúdo de choque. Em um dos registros de referência o eixo $x$ dá 146
contra 6,3 no eixo $z$ --- o sinal longitudinal é dominado por eventos
transientes raros e severos, exatamente a situação que a soma de sexta potência
foi concebida para capturar e que uma análise de valor eficaz trataria como
equivalente a um sinal muito mais benigno.

\paragraph{Tela cheia.} F11, duplo clique no gráfico, ou o botão \emph{Full
screen} na barra de navegação alternam qualquer gráfico para uma janela
dedicada, dimensionada à tela e com margens fixas; Esc, F11, ou o botão de
fechar no canto direito da barra de navegação, adicionado nesta versão, voltam.
A figura é
o mesmo objeto do \code{matplotlib} nos dois modos, de modo que zoom e
deslocamento aplicados sobrevivem à transição --- o que exigiu, nesta versão,
trocar o motor de layout de restrito (\emph{constrained}) para margens fixas
explícitas apenas durante a tela cheia, porque o motor de layout automático
recalcula margens de forma inconsistente quando a figura muda de tamanho e de
container, chegando a cortar o rótulo do eixo horizontal em gráficos de
frequência.

\paragraph{Barra de progresso.} Durante qualquer operação em segundo plano
(Seção~\ref{sec:threading}), a barra de status mostra a etapa em andamento ---
qual ponto de contato, qual arquivo, qual porcentagem do cálculo NHO-09 ---
o que importa em especial para registros longos.

\paragraph{Enquadramento de gráfico persistente.} Os botões de \emph{pan} e
\emph{zoom} mostram o estado ligado em azul, e o enquadramento sobrevive ao
redesenho: mudar um controle redesenha a figura, o que antes descartava
qualquer zoom aplicado pelo usuário. O enquadramento é agora guardado e
reposto a cada redesenho; o botão \emph{Home} volta ao completo.

\paragraph{Tema.} As cores ficam em \code{theme.py}, compartilhadas pelos dois
executáveis. A paleta é clara e explícita de propósito: sem isso o Qt herda o
tema do sistema operacional, e no Windows em modo escuro o texto claro do tema
cai em cima dos fundos claros da folha de estilo e some.

% =====================================================================
% =====================================================================
\section{Lições de estudos semelhantes}
\label{sec:licoes}

O grupo mantém um acervo de artigos sobre vibração de corpo inteiro, reunido
principalmente para orientar a linha de pesquisa pediátrica
(Seção~\ref{sec:criancas}), mas que também documenta, em estudos de campo
reais publicados por outros grupos, os mesmos problemas metodológicos que este
software enfrentou de forma independente. Esta seção extrai, artigo por
artigo, a lição relevante para uma decisão concreta de projeto do software ---
não um resumo genérico da literatura.

\subsection{Extrapolação de amostra parcial para dose diária: precedente publicado}

A Fórmula~(5) da Cláusula~5 (Seção~\ref{sec:cadeiaC5}) extrapola uma dose
medida por um tempo $t_m$ para um tempo de exposição diário $t_d$ --- é
exatamente o mecanismo que a importação de campo usa para tratar os blocos de
\unit{2}{s} do registrador como uma condição medida (item concluído em
\emph{Trabalho futuro}). Três estudos independentes usam a mesma lógica sobre
amostras de campo tão ou mais fragmentadas que as nossas: de la
Hoz-Torres et al. (2022) combinam passagens repetidas de \unit{10{,}5}{km} por
superfície via a mesma razão $t_d/t_m$ da norma; Singh et al. (2022) gravam
apenas \unit{60}{s} por ensaio ``devido à disponibilidade limitada de campo'' e
alimentam o resultado direto na Cláusula~5/Anexo~C; Rao, Sivapirakasam e Arora
(2018) extrapolam trajetos parciais de ônibus escolar para a rotina diária da
criança pela fórmula de VDV estimado, $\mathrm{eVDV} = 1{,}4\,a_v\,T^{1/4}$. Em
nenhum dos três a extrapolação em si é questionada pelos autores ou por
revisores --- o que é questionado, sempre, é se a janela medida representa bem
a atividade. Isso não resolve a questão em aberto sobre a representatividade
estatística dos 40\% de blocos mantidos pelo nosso registrador
(Seção~\ref{sec:campo}), mas mostra que a extrapolação pela Fórmula~(5) em si
segue prática já estabelecida, e não é uma aposta metodológica isolada deste
software.

\subsection{Seleção automática de caminho: um ponto cego já documentado na literatura}
\label{sec:pontocego}

O critério que decide entre o Anexo~A/E e a Cláusula~5/Anexo~C
(Seção~\ref{sec:validade}) e a checagem de validade deste software
(\code{core.validity\_check}, Seção~\ref{sec:validade}) dependem exclusivamente
do pico ponderado do eixo $z$. De la Hoz-Torres et al. (2019), comparando a
mesma gravação de um trator agrícola pelas edições de 2004 e 2018 da
ISO~2631-5, relatam um detalhe que confirma isso como um ponto cego real da
norma, não uma simplificação de implementação: no terreno irregular medido, os
eixos horizontais superaram o vertical ($a_x = 1{,}18$, $a_y = 0{,}97$ contra
$a_z = 0{,}85\,\mathrm{m/s^2}$ eficazes), mas o critério de regime --- severo ou
menos severo --- da edição de 2018 olha só para o pico de $z$. Um evento
dominado por $x$ ou $y$ pode passar despercebido pelo mesmo teste que decide
qual cadeia de cálculo rodar. Este software herda esse ponto cego porque
implementa o critério da norma tal como escrito; a Seção~\ref{sec:zonly}
registra isso como questão em aberto.

\subsection{Diagnóstico em frequência ao lado de um método no domínio do tempo}

A Cláusula~5/Anexo~C é deliberadamente um método no domínio do tempo: nenhuma
ponderação em frequência é aplicada (docstring de \code{annexc.py}). Dois
estudos mostram que essa escolha, apesar de correta para o cálculo da dose,
pode esconder fenômenos que só aparecem olhando o espectro. Delcor et al.
(2022), avaliando desconforto de passageiros de helicóptero, mostram que o
índice de conforto da ISO~2631-1 --- puramente espectral --- subestima
sistematicamente o desconforto quando há batimento (modulação de amplitude
entre a frequência do rotor principal e do rotor de cauda); corrigir isso
exigiu sobrepor uma camada de detecção de envelope no domínio do tempo,
elevando o $R^2$ de $0{,}78$ para $0{,}84$. Ittianuwat, Fard e Kato (2016)
mostram o problema espelhado: o ponto único de medição do encosto prescrito
pela ISO~10326-1 cai perto de um nó estrutural da ressonância de torção do
banco, subestimando a vibração ali medida em 94\% a 229\% frente a outros
pontos próximos do mesmo encosto, nas frequências de \unit{31}{a}
\unit{37{,}5}{Hz}. As duas lições apontam na mesma direção: um método de
domínio único e ponto único pode perder fenômenos físicos reais. Foi essa
literatura que motivou o novo gráfico \emph{Frequency spectrum} na aba
\emph{Clause 5 / Annex C} desta versão --- puramente informativo, calculado
sobre o $a_z$ bruto (não condicionado) via \code{core.welch\_spectrum}, sem
alimentar nenhum resultado da Cláusula~5. Para os dados de campo, cada bloco de
\unit{2}{s} entra como um segmento de Welch próprio, o mesmo padrão já usado no
painel NHO-09 (Seção~\ref{sec:fieldlogger}).

\subsection{Relatar mais de uma métrica, não só uma}

De la Hoz-Torres et al. (2022) computam $A(8)$, VDV e $S_d^A/R^A$ sobre a
mesma gravação e obtêm vereditos diferentes: VDV foi a métrica mais
conservadora, sinalizando risco em superfícies onde $A(8)$ e $S_d^A$ ficaram
dentro da faixa segura. Os autores reconhecem que isso ``pode causar confusão
ao gestor de segurança'', mas optam por reportar as três em vez de escolher
uma. Este software já segue essa prática no painel de campo, mostrando
NHO-09 e Diretiva~2002/44/CE lado a lado (Seção~\ref{sec:fieldlogger}); o novo
gráfico \emph{Risk reference bands} da aba \emph{Clause 5 / Annex C}
(Seção~\ref{sec:contato}) estende o mesmo princípio para dentro de uma única
métrica: em vez de mostrar só $\Pi$ como um número isolado, mostra onde o $R$
calculado cai entre os três pontos de calibração da Tabela~C.2
($\Pi = 10\%/50\%/90\%$), tornando explícito que o limiar de ``baixo risco''
não é um número redondo arbitrário, mas um ponto específico da própria tabela
da norma.

\subsection{Limiares concretos de dano: o que já existe publicado}

Godley e Csongradi (2023), revisando terapia por vibração de corpo inteiro em
crianças com deficiência, recalculam a intensidade (em $g$) de protocolos
terapêuticos publicados a partir da frequência e do deslocamento do
equipamento, já que a maioria dos estudos nunca reporta a intensidade
diretamente. O limite convencional da ISO~2631 (\unit{1}{g}) serve de
referência: um protocolo documentado como seguro (Mogil et al., estudo VIBE,
65 crianças, um ano, sem eventos adversos, com ganho de densidade óssea) opera
a \unit{0{,}3}{a}~\unit{0{,}4}{g}, enquanto vários protocolos publicados
chegam a \unit{15}{a}~\unit{240}{g} --- perigosamente acima. Os autores também
citam faixas de ressonância associadas a dano (\unit{\sim 5}{Hz} no tronco,
abaixo de \unit{20}{Hz} nos membros inferiores, acima de \unit{50}{Hz} com
risco de dano muscular) e efeitos adversos documentados em magnitudes altas
(deslocamento de disco intervertebral, degeneração espinhal, ruptura de
retina, hemorragia intracerebral). Esses números são do contexto de terapia
pediátrica, numa escala diferente da ocupacional adulta que este software
calcula, e não substituem a Tabela~C.2 --- mas são exatamente o tipo de
``o que é perigoso'', com número e fonte, que faltava ao lado do resultado do
assessment. Ficam registrados aqui como o conjunto de referência natural para
quando o \emph{Modelo pediátrico experimental} (\emph{Trabalho futuro}) existir
e precisar de seus próprios limiares, em vez dos da Tabela~C.2 adulta.

\subsection{Parâmetros biomecânicos pediátricos: sempre escalados, nunca medidos diretamente}

Dois modelos de elementos finitos da coluna pediátrica/adolescente, lidos para
avaliar se existe algum parâmetro biomecânico infantil publicado que pudesse
alimentar uma extensão do Anexo~C (Seção~\ref{sec:criancas}), confirmam de
forma independente a mesma limitação. Muhayudin et al. (2021) constroem um
modelo da coluna torácica de uma criança de 8--9 anos a partir de um modelo
físico Sawbones, mas as propriedades de material do osso, ligamentos e disco
não foram medidas: foram fixadas como frações do valor adulto (osso
trabecular/cortical a 57--66\%, ligamentos a 90\%) ``porque os dados da
literatura eram limitados''. Aroeira et al. (2018, UFMG), modelando a coluna
torácica de uma adolescente de 10 anos sob hipocifose, escalam a geometria
vertebral a 71--91\% das dimensões adultas a partir de uma única referência
(Meijer et al.) e afirmam textualmente que ``há pouca informação biomecânica
sobre colunas vertebrais jovens''. As duas fontes, de grupos independentes,
chegam à mesma prática: nenhum parâmetro biomecânico pediátrico da coluna é
medido diretamente na literatura disponível --- todos são escalados de dados
adultos por uma razão fixa. Qualquer extensão pediátrica deste software
(Seção~\ref{sec:criancas}) herdaria exatamente essa limitação, e deve
declará-la, em vez de implicar que existe uma calibração pediátrica validada.
Separadamente, Zhang et al. (2026), modelando o corpo sentado adulto por
dinâmica multicorpos, mostram que calibrar um modelo segmentado só contra uma
métrica de corpo inteiro (massa aparente, transmissibilidade assento-cabeça)
deixa os parâmetros internos por segmento não unicamente determinados, a menos
que uma medição local independente seja usada para desambiguar --- um
argumento a mais para que a sincronização multiponto do registrador de campo
(arquitetura master-slave proposta no artigo CONEM2026-0596, Seção
\ref{sec:conjuntoB}) valha não só pela cobertura espacial do Anexo~A, mas
para validar as próprias premissas do modelo de ponto único do Anexo~C.

\subsection{Nenhum limiar pediátrico validado existe, em nenhuma direção}

Seis fontes independentes, lidas nesta revisão e na que originou a
Seção~\ref{sec:criancas}, convergem no mesmo diagnóstico a partir de ângulos
opostos. Pelo lado do risco: Zuska et al. (2021) medem acelerações em
cadeirinhas infantis reais que excedem o limite adulto de conforto da
ISO~2631-1 (\unit{0{,}315}{m/s^2}) em quase toda condição de pista testada, e
afirmam explicitamente que ``não existem normas assim para crianças''; Bressel,
Smith e Branscomb (2010) e Rao, Sivapirakasam e Arora (2018), já citados na
Seção~\ref{sec:criancas}, mostram transmissibilidade e exposição maiores em
crianças mais novas. Pelo lado do benefício: Pulay et al. (2023), em uma
meta-análise de 16 ensaios clínicos de terapia por vibração em crianças com
paralisia cerebral, confirmam ganho funcional mensurável em doses baixas, mas
constatam que a exposição nunca foi quantificada de forma rigorosa nos ensaios
revisados (heterogeneidade $I^2$ de até 96\%). Aplicar o limiar adulto a
crianças, portanto, tanto superestima quanto subestima o risco real dependendo
do contexto --- e nenhuma das duas direções tem lastro normativo. Esse é o
argumento mais forte, entre tudo o que esta revisão levantou, para que
qualquer saída pediátrica deste software seja rotulada como estimativa do
laboratório, nunca como aplicação certificada da norma, exatamente como a
Seção~\ref{sec:criancas} já recomenda.

\subsection{Área de placa terminal: também não é previsível por antropometria}

Seidel, Pöpplau, Morlock, Püschel e Huber (2008) tentam exatamente o que uma
extensão pediátrica do Anexo C precisaria fazer --- prever a área de placa
terminal ($B$ na Fórmula~C.1, \code{endplate\_area} neste software) a partir de
características antropométricas externas --- em 53 colunas de cadáveres
masculinos (\unit{33{,}2 \pm 5{,}8}{anos}), medidas por tomografia
computadorizada. O resultado é negativo: nenhuma correlação simples ou múltipla
entre massa corporal, altura, ou diâmetros de punho/cotovelo/joelho/tornozelo e
a área de placa explica mais que 20\% da variância, e o homem mais leve da
amostra (\unit{52{,}3}{kg}) tinha uma placa de L4 \emph{maior} que o mais pesado
(\unit{135}{kg}). Os autores concluem que características antropométricas
externas não servem para prever o tamanho da placa terminal individual. Isso
não é um problema deste software --- é uma propriedade real da anatomia
humana --- mas significa que os dois valores fixos de \code{endplate\_area}
(14~cm² homem, 11,6~cm² mulher, Fórmula~C.2) não têm como ser refinados por
massa ou altura mesmo se quiséssemos, e que qualquer tentativa de escalar a
área de placa por antropometria numa extensão pediátrica enfrentaria a mesma
ausência de correlação nos poucos dados adultos onde isso já foi testado.

\subsection{A Tabela C.2 rastreada até a fonte}

Schmidt, Paskoff, Shender e Bass (2012, \emph{Spine} 37:E1614--E1621) --- a
referência [66] da norma --- são a fonte exata da Tabela~C.1/C.2 e da
Fórmula~C.4 deste software: os coeficientes $\alpha$/$\beta$ de Weibull
(1,613/2,799 homens; 0,959/3,709 mulheres) e os valores de $R$ em
10\%/50\%/90\% (0,72/1,42/2,17 homens; 0,52/0,87/1,20 mulheres) publicados na
Tabela~5 do artigo batem exatamente com os que este software já reproduzia
(Seção~\ref{sec:auditoria}) --- confirmando a rastreabilidade completa, não só
a reprodução numérica. A amostra é 105 segmentos vertebrais de 77 homens e 28
mulheres, 19 a 93 anos --- a mesma faixa etária que a Seção~\ref{sec:criancas}
já identifica como excluindo crianças e extrapolando na direção errada para
elas ($S_u$ maior, risco aparente menor).

\subsection{Validação epidemiológica real: os limiares de R podem não ser protetores}

Bovenzi, Schust, Menzel, Prodi e Mauro (2015, \emph{Int Arch Occup Environ
Health} 88:487--499) acompanham 537 motoristas profissionais por dois anos e
comparam a incidência real de dor lombar contra $R$ (calculado pelo mesmo
modelo do Anexo C) e contra $A(8)$/VDV da Diretiva. A dose interna ($S_{ed}$,
$R$) prevê melhor os desfechos de dor lombar do que as métricas externas da
Diretiva --- um resultado a favor da abordagem biomecânica que este software
implementa. Mas o achado mais importante é outro: o quartil superior de $R$
observado na coorte (0,41--0,72) já mostra o dobro da incidência de dor lombar
em relação ao quartil inferior (raz\~ao de chances $\geq 2{,}0$, $p<0{,}05$) ---
e esse quartil superior está \emph{inteiro abaixo} do limiar de baixo risco que
a norma usa, $R<0{,}8$. Os próprios autores concluem: ``the current boundary
values for the risk factor R proposed by ISO/CD 2631-5 (2014) are not
protective for the health of the lumbar spine of the exposed workers''. Schust
et al. (2015, \emph{Ergonomics} 58:1191--1206), reanalisando a mesma coorte com
o modelo de elementos finitos completo do Anexo A/E em vez da fórmula
simplificada, confirmam a mesma ordem de importância dos fatores (aceleração,
duração da exposição, postura, idade de início e massa/IMC, nessa ordem) e a
mesma associação fraca entre $R$ e $A(8)$/VDV. Para este software, que já
mostra $R$ e $\Pi$ ao lado dos limiares 10\%/50\% na aba \emph{Clause 5 / Annex
C} (Seção~\ref{sec:licoes}), esse achado é um argumento para nunca apresentar
$\Pi<10\%$ como ``seguro'' sem qualificação --- é o que a própria norma
permite chamar de baixa probabilidade, mas a única validação epidemiológica
publicada contra essa fronteira sugere que ela já não é conservadora o
bastante.

\subsection{Duas confirmações independentes, uma de modelo e uma \emph{in vivo}}

Fleury (2015, 50\textsuperscript{a} UK Conference on Human Responses to
Vibration) reproduz, com um método algébrico alternativo, os mesmos resultados
do software de referência da DIN SPEC 45697 --- a norma alemã da qual o
Anexo A herda o modelo --- e faz um estudo paramétrico do risco por postura,
antropometria, frequência, direção e ponto de aplicação da excitação. Duas
coisas confirmam achados já registrados aqui de forma independente: a Figura~2
do artigo, que traça os limites de massa por classe de IMC do DIN SPEC 45697,
bate visualmente com a Tabela~2 de Seidel et al. (2008) usada para resolver o
índice antropométrico (Seção~\ref{sec:indice}); e o efeito de postura é
quantificado --- a \unit{3{,}5}{Hz}, um motorista no grupo~1 leva 29 anos para
atingir $R=1{,}2$, contra 22 anos no grupo~3, mesma aceleração.

Rohlmann, Hinz, Blüthner, Graichen e Bergmann (2010, \emph{European Spine
Journal} 19:1129--1135) --- referência [53] da norma --- medem isso
\emph{de verdade}: um implante de substituição vertebral telemetrizado, em
quatro pacientes reais, mede a força de compressão durante vibração controlada
em mesa hexápode. Encostar no encosto reduz a carga no implante para 48--60\%
da que ocorre sentado sem encosto; vibração multieixo aumenta a carga em
23--89\% frente a um único eixo, na mesma ordem de grandeza da soma multiponto
da Fórmula~A.1. Não é uma correção para este software --- é a única fonte,
entre tudo o que foi revisado, com uma medida direta em corpo humano vivo, e
confirma que os efeitos de postura e multieixo que o Anexo A já modela por
elementos finitos são reais, não artefatos do modelo.

% =====================================================================
\section{Questões em aberto}
\label{sec:limitacoes}

\subsection{O critério de validade considera só o eixo z}
\label{sec:zonly}

O teste que decide entre o Anexo~A/E e a Cláusula~5/Anexo~C
(\code{core.validity\_check}, Seção~\ref{sec:validade}) olha exclusivamente
para o pico ponderado do eixo $z$, porque é assim que a Seção~6.2 da norma
define o critério. A Seção~\ref{sec:pontocego} registra um caso publicado --- de
la Hoz-Torres et al. (2019), comparando as edições de 2004 e 2018 da
ISO~2631-5 sobre a mesma gravação de trator agrícola --- em que os eixos
horizontais superaram o vertical, e o mesmo tipo de critério (baseado só em
$z$) não teria detectado isso. Este software implementa a norma tal como
escrita, então herda o mesmo ponto cego: um evento dominado por $x$ ou $y$ ---
plausível em terreno irregular, e não descartável nos blocos de \unit{2}{s} do
registrador de campo (Seção~\ref{sec:campo}) --- pode ficar abaixo do limiar de
\unit{9{,}81}{m/s^2} em $z$ mesmo excedendo-o largamente em outro eixo, sem que
o software sinalize nada de anormal. Não há correção proposta aqui, porque
alterar o critério divergiria do texto da norma; fica registrado como
limitação conhecida, e como argumento a mais para inspecionar visualmente os
três eixos (por exemplo pelo gráfico de espectro em frequência recém-adicionado
à aba \emph{Clause 5 / Annex C}, Seção~\ref{sec:licoes}) em vez de confiar cegamente
no resultado do \code{validity\_check}.

\subsection{A decomposição do índice antropométrico}
\label{sec:indice}

\textbf{Esta seção substitui uma versão anterior desta documentação, que
concluía que a decomposição do índice não estava determinada. Uma fonte
publicada, encontrada depois, resolve a questão.}

Propôs-se em algum momento que o índice $k$ de 30 posições se decompusesse como
$15\,(\mathrm{classe\ de\ IMC} - 1) + 3\,(\mathrm{percentil} - 1) +
\mathrm{subclasse}$, com duas classes de índice de massa corporal separadas pelo
limiar de \unit{26{,}1}{kg/m^2}. Testada por busca exaustiva contra os únicos
casos de exemplo que variam a massa corporal (\code{indiv5} e \code{indiv6}, que
percorrem 70, 80, 90, 100 e \unit{110}{kg} a \unit{1{,}85}{m}), essa hipótese
parecia quase certa mas não exata: o melhor ajuste reproduzia os dois casos com
erro relativo de $6\times10^{-4}$ --- vinte vezes maior que os $3\times10^{-5}$
dos casos que fecham exatamente --- e colocava \emph{todas} as massas no
segundo bloco, inclusive \unit{70}{kg} a \unit{1{,}85}{m} (índice de massa
corporal de 20,5), o que não fazia sentido se o bloco representasse classe de
IMC. Como os oito arquivos de exemplo disponíveis usam sempre a mesma altura,
massa e IMC ficam perfeitamente confundidos entre si, e nenhuma busca adicional
sobre esses dados poderia separar as duas hipóteses.

O texto da norma (Anexo A.3.1 a) confirma o limiar de \unit{26{,}1}{kg/m^2}
diretamente, mas não publica a tabela completa de percentis. A referência
[57] da própria norma --- Seidel, Hinz, Hofmann, Menzel (2008), \emph{Intraspinal
forces and health risk caused by whole-body vibration}, Int J Ind Ergon
38:856--867 --- publica essa tabela (sua Tabela 2), derivada dos 1.163
motoristas europeus que originaram as funções de transferência do Anexo A. A
peça que faltava: o percentil de massa não é calculado numa escala única e
global, mas \emph{dentro} da classe de IMC já escolhida --- as duas escalas se
sobrepõem (o 5\textsuperscript{o} percentil da classe alta, \unit{78}{kg}, fica
abaixo do 75\textsuperscript{o} percentil da classe baixa, \unit{82}{kg}).
Substituindo o ajuste
empírico por essa tabela, \unit{70}{kg} a \unit{1{,}85}{m} cai corretamente na
classe baixa de IMC, sem forçar nada.

A correção foi validada contra os \emph{seis} casos individuais da norma
(\code{indiv1} a \code{indiv6}, não só os dois que variam massa), rodando o
histórico completo de cada um --- antropometria e padrão de exposição variando
ano a ano por 45 anos --- pelo mesmo \code{analysis.Runner} que a interface usa,
e comparando o $R^A$ obtido contra o valor oficial publicado em
\code{Results\_ISO.xls}. Os seis batem com erro relativo entre $7\times10^{-6}$
e $3{,}5\times10^{-5}$ --- o mesmo patamar de precisão numérica dos casos já
validados na Seção~\ref{sec:auditoria}, não mais uma aproximação. O script está
em \code{tests/anthropometry\_check.py}.

A análise de sensibilidade da Seção~\ref{sec:sensibilidade} continua sendo o
resultado mais honesto para populações fora do padrão de calibração do modelo
--- usuários de cadeira de rodas, crianças --- porque nenhuma quantidade de
dados resolve o que fazer quando o sujeito real não se parece com nenhum dos
30 modelos. Mas deixou de ser necessária como paliativo para uma incerteza na
própria decomposição do índice: essa incerteza não existe mais.

\subsection{A terceira dimensão do índice}

A hipótese mais provável é classe de magnitude de aceleração: a norma menciona
que as funções de transferência foram calculadas para diferentes classes de
magnitude, atribuídas automaticamente pelo software a partir das acelerações
medidas, e o código-fonte original contém uma função que classifica o sinal em
três faixas segundo o valor eficaz, com limiares em \unit{0{,}65}{} e
\unit{1{,}35}{m/s^2}.

Entretanto, \emph{todos} os casos de referência coincidem com a subclasse 1,
inclusive o do bulldozer, cujo valor eficaz de \unit{3{,}08}{m/s^2} deveria
situá-lo na terceira faixa. A implementação adota a subclasse 1 como padrão e
expõe a escolha na interface.

\subsection{Duração de medição}

A norma define $t_m$ como a duração da medição. O software original, contudo,
usa o valor declarado no campo correspondente da planilha, ainda que não
coincida com o comprimento efetivo do registro: nos exemplos de referência
declara-se \unit{400}{s} para registros de \unit{140}{s}. O comportamento foi
mantido por compatibilidade, e constitui ponto de atenção na preparação de dados
próprios.

\subsection{Proteção da população}

O campo correspondente na interface \emph{não decorre da norma}. A ISO 2631-5
menciona a proteção de 50\% da população apenas como característica embutida nos
coeficientes da Fórmula (E.2), sem fornecer formulação para outros percentis. O
campo existe porque o software original o expunha e o registrava nos arquivos de
saída. Em 50\% o resultado é exatamente o da norma; em outros valores aplica-se
uma correção normal-padrão que não pôde ser validada contra referência.

\subsection{Aplicação a crianças}
\label{sec:criancas}

Uma linha de pesquisa em andamento no grupo pretende avaliar exposição em
crianças, inclusive usuárias de cadeira de rodas. Ambos os caminhos da norma têm
base biomecânica adulta, e isso aparece de forma concreta na
Equação~\eqref{eq:c3}: a relação foi derivada de ensaios de fadiga em espécimes
cadavéricos de 19 a 93 anos. Para uma criança de 10 anos, $b+i = 10$ está fora do
intervalo dos dados, e a extrapolação devolve resistência \emph{maior}, fazendo
o risco parecer \emph{menor} --- a direção errada de errar. Vértebras infantis
têm o anel apofisário ainda não fundido, um ponto de fragilidade conhecido sob
carga axial. O modelo de Weibull da Equação~\eqref{eq:weibull} tem a mesma
origem.

O caminho do Anexo C é o menos inadequado, porque o $m_z$ é parâmetro aberto e
admite antropometria arbitrária. Ainda assim, nenhum resultado nesse domínio
deveria ser apresentado como aplicação da norma sem discussão explícita desta
extrapolação.

Um levantamento de literatura conduzido para orientar esta linha de pesquisa
confirma e refina esse diagnóstico; a Seção~\ref{sec:licoes} trata cada fonte em
detalhe, esta seção resume as conclusões. Múltiplos estudos independentes ---
crianças em pé sobre plataforma vibratória (Bressel, Smith e Branscomb, 2010),
crianças sentadas em ônibus escolar (Rao, Sivapirakasam e Arora, 2018),
crianças em cadeirinhas automotivas (Zuska et al., 2021) --- mostram que a
transmissibilidade corpo-assento em crianças difere sistematicamente da de
adultos, tipicamente por menor massa de tecido amortecedor: a transmissibilidade
tronco/quadril pode superar em 42--62\% a de adultos em certas bandas de
frequência (Bressel et al.), e a exposição medida é sistematicamente maior em
crianças mais novas e mais leves (Rao et al.) --- embora nem todo segmento
corporal mostre essa diferença de forma uniforme. Publicações que citam
explicitamente a ISO 2631-1 e a ISO 2631-5 afirmam, sem exceção, que nenhuma
delas foi validada para população pediátrica, e aplicam as tabelas adultas
apenas por falta de alternativa, sempre com ressalva textual nesse sentido.

Dessa revisão emerge um caminho técnico, ainda não implementado, que aproveita
a própria arquitetura descrita na Seção~\ref{sec:dados}: a cadeia do Anexo A não
tem nenhuma premissa de adulto embutida no código --- ela consome um conjunto de
funções de transferência no formato \code{BSKUefuM\_k\_p\_d}
(Seção~7.1) e faz o resto. Um conjunto equivalente, derivado de um modelo
biodinâmico pediátrico calibrado por antropometria (massa, estatura, idade) em
vez de calibrado por sujeito individual, alimentaria \code{core.py} sem
alteração na lógica de cálculo. Um modelo multicorpos compacto publicado
recentemente para o corpo adulto sentado (Zhang et al., 2026) já propõe
justamente esse tipo de mapeamento antropometria$\to$parâmetro, e conjuntos
independentes de propriedades mecânicas da coluna torácica pediátrica e
adolescente (módulo de elasticidade do osso cortical, trabecular, disco e
ligamentos) já foram publicados a partir de modelos de elementos finitos
(Muhayudin et al., 2021; Aroeira et al., 2018) --- os mesmos dois ingredientes,
massa-mola-amortecedor e elementos finitos, que originaram o modelo adulto do
Anexo A. Nenhuma dessas duas fontes de elementos finitos mede a propriedade
pediátrica diretamente: ambas escalam dado adulto por uma razão fixa e dizem
isso com todas as letras (Seção~\ref{sec:licoes}) --- uma extensão deste
software herdaria a mesma limitação e deve declará-la. Qualquer resultado
produzido por essa via deveria ser explicitamente rotulado como estimativa do
laboratório, não como aplicação da norma, do mesmo modo como a Seção~11.4 trata
o campo de proteção populacional: uma extensão consistente com o espírito do
método, mas sem o mesmo lastro normativo.

\subsection{Qualidade do material de origem}

A norma e o código-fonte foram consultados a partir de extração por
reconhecimento óptico. O corpo do texto e as fórmulas apresentam boa qualidade,
mas diversas tabelas e trechos de código foram degradados --- a inconsistência
do $W_d$ a \unit{10}{Hz} é provavelmente um caso. Valores numéricos citados
devem ser conferidos no material original antes de uso em publicação.

% =====================================================================
\section{Trabalho futuro}

Quatro dos cinco itens de trabalho futuro listados na versão anterior desta
documentação foram concluídos e estão descritos nas seções indicadas; o item
restante, e os que se abriram desde então, seguem em ordem de prioridade
sugerida:

\begin{enumerate}
  \item \textbf{[Concluído]} Integração com a instrumentação do laboratório ---
        leitura direta do formato \code{.m} gravado pelo registrador
        (Seções~\ref{sec:campo} a~\ref{sec:fieldlogger}). A sincronização entre
        os três sensores permanece limitada por hardware
        (Seção~\ref{sec:conjuntoB}) e é a fonte mais provável de rejeição por
        deriva de base de tempo quando o Anexo A é aplicável.
  \item \textbf{[Concluído]} Análise de sensibilidade antropométrica e postural
        (Seção~\ref{sec:sensibilidade}), executando o mesmo sinal contra os 30
        modelos e as sete posturas e quantificando a dispersão do $R^A$.
  \item \textbf{[Concluído]} Resolver o índice antropométrico
        (Seção~\ref{sec:indice}). A Tabela~2 de Seidel, Hinz, Hofmann, Menzel
        (2008) --- referência [57] da própria norma --- publica os percentis
        de massa corporal dentro de cada classe de IMC; substituindo o ajuste
        empírico anterior por essa tabela, os seis casos individuais da norma
        (\code{indiv1}-\code{indiv6}) reproduzem o $R^A$ oficial com erro
        relativo de $10^{-5}$-$10^{-6}$, o mesmo patamar dos casos já
        validados. Não foi preciso contatar a instituição original.
  \item \textbf{Modelo pediátrico experimental.} Construir e validar, como
        módulo paralelo e claramente rotulado, um conjunto de funções de
        transferência análogo ao do Anexo A mas calibrado para antropometria
        infantil (Seção~\ref{sec:criancas}), a partir de um modelo
        multicorpos parametrizável e de propriedades mecânicas da coluna
        pediátrica já publicadas. É o item de maior potencial de contribuição
        original do grupo, dado que a lacuna está bem documentada na literatura
        mas nenhuma fonte revisada até agora a preenche.
  \item \textbf{[Concluído]} Importação direta de \code{.m} pela Cláusula
        5/Anexo C. Implementada na aba \emph{Clause 5 / Annex C} (não no
        construtor de planilhas, por decisão deliberada --- ver
        Seção~\ref{sec:licoes}): cada bloco de \unit{2}{s} do registrador de
        campo é tratado como uma mini-exposição medida, combinada pela mesma
        extrapolação $t_d/t_m$ (Fórmula~5) que a norma já usa entre exposições
        distintas. O estudo controlado da Seção~\ref{sec:campo} mostrou que o
        reinício do filtro passa-alta a cada bloco desvia menos de 1\% --- não
        é obstáculo --- e a revisão de literatura da Seção~\ref{sec:licoes}
        encontrou precedente publicado para exatamente esse tipo de
        extrapolação de amostra parcial (60\,s, poucos minutos de trecho, ou um
        ciclo de trabalho de 40\%) para uma dose diária, usando a mesma lógica
        de escala temporal da própria norma. Um teste isolado em dados reais de
        campo (dia inteiro, assento, 25/03/2025) produziu $R=0{,}59$,
        $\Pi\approx 6\%$ --- dentro da faixa plausível da Tabela~C.2, nem
        saturado em zero nem em cem por cento --- antes de a funcionalidade ser
        exposta na interface.
  \item \textbf{Processamento segmentado com sobreposição e soma}, se houver
        interesse em janelar sem perder a reconstrução no tempo, com
        revalidação (Seção~\ref{sec:espectral}).
  \item \textbf{Instrumentação de detecção de perda de contato}, substituindo a
        heurística da Seção~\ref{sec:contato} por medida direta.
\end{enumerate}

\vspace{1em}
\noindent\textbf{Referências normativas.}\;
ISO 2631-5:2018 --- \emph{Mechanical vibration and shock: Evaluation of human
exposure to whole-body vibration. Part 5: Method for evaluation of vibration
containing multiple shocks}. ISO, 2018.\\
ISO 2631-1:1997 --- \emph{Mechanical vibration and shock: Evaluation of human
exposure to whole-body vibration. Part 1: General requirements}. ISO, 1997.\\
NHO-09 --- \emph{Avaliação da Exposição Ocupacional a Vibrações de Corpo
Inteiro}. Fundacentro, 2013.\\
Diretiva 2002/44/CE do Parlamento Europeu e do Conselho, de 25 de junho de 2002,
relativa às prescrições mínimas de segurança e saúde respeitantes à exposição
dos trabalhadores aos riscos devidos a agentes físicos (vibrações).

\vspace{1em}
\noindent\textbf{Artigos revisados (Seções~\ref{sec:indice} e~\ref{sec:licoes}).}\;
\begin{enumerate}\itemsep2pt
\item Avila, J.P.R.; Lima~II, E.J.; Donadon, L.V.; Duarte, M.L.M.; Augusto,
      S.O.; Huang, F.; Pereira, P.T.S. \emph{Development of a Low-Cost
      Alternative for Measuring Human Whole-Body Vibration (WBV): Problems and
      Perspectives}. CONEM2026-0596, 13\textsuperscript{o} Congresso Nacional
      de Engenharia Mecânica, 2026.
\item Aroeira, R.M.C.; Pertence, A.E.M.; Kemmoku, D.T.; Greco, M. \emph{The
      Effect of Hypokyphosis on the Biomechanical Behavior of the Adolescent
      Thoracic Spine}. Journal of the Brazilian Society of Mechanical Sciences
      and Engineering, 2018.
\item Bovenzi, M.; Schust, M.; Menzel, G.; Prodi, A.; Mauro, M.
      \emph{Relationships of low back outcomes to internal spinal load: a
      prospective cohort study of professional drivers}. International
      Archives of Occupational and Environmental Health, 88, 487--499, 2015.
      --- referência [69] da ISO 2631-5:2018; conclui que os limiares de $R$
      da norma podem não ser protetores (Seção~\ref{sec:licoes}).
\item Bressel, E.; Smith, G.; Branscomb, J. \emph{Transmission of Whole Body
      Vibration in Children While Standing}. Clinical Biomechanics, 25,
      181--186, 2010.
\item Fleury, G. \emph{Parametric study of the spinal load model used in the
      draft standard ISO2631-5}. Proceedings of the 50\textsuperscript{th}
      UK Conference on Human Responses to Vibration, University of
      Southampton, 2015. --- referência [76] da ISO 2631-5:2018; estudo
      paramétrico do mesmo modelo do Anexo~A (Seção~\ref{sec:licoes}).
\item de la Hoz-Torres, C.; Aguilar-Aguilera, A.J.; Martínez-Aires, M.D.;
      Ruiz, D.P. \emph{A Comparison of ISO 2631-5:2004 and ISO 2631-5:2018
      Standards for Whole-Body Vibrations Exposure: A Case Study}. In:
      Occupational and Environmental Safety and Health, Studies in Systems,
      Decision and Control, vol.~202. Springer, 2019.
\item de la Hoz-Torres, C.; Aguilar, A.J.; Ruiz, D.P.; Martínez-Aires, M.D.
      \emph{Whole Body Vibration Exposure Transmitted to Drivers of Heavy
      Equipment Vehicles: A Comparative Case According to the Short- and
      Long-Term Exposure Assessment Methodologies Defined in ISO 2631-1 and
      ISO 2631-5}. International Journal of Environmental Research and Public
      Health, 19, 5206, 2022.
\item Delcor, L.; Parizet, E.; Ganivet-Ouzeneau, M.; Caillet, D.
      \emph{Helicopter vibrations: study on the use of standard ISO 2631-1 to
      estimate discomfort}. Ergonomics, 65(2), 296--304, 2022.
\item Godley, A.; Csongradi, J. \emph{Whole Body Vibration Therapy for
      Children with Disabilities: A Survey of Potential Risks and Benefits}.
      Archives of Rehabilitation Research and Clinical Translation, 2023.
\item Hinz, B.; Seidel, H.; Hofmann, J.; Menzel, G. \emph{The significance of
      using anthropometric parameters and postures of European drivers as a
      database for finite-element models when calculating spinal forces
      during whole-body vibration exposure}. International Journal of
      Industrial Ergonomics, 38, 816--843, 2008. --- referência [27] da ISO
      2631-5:2018; fonte do limiar de IMC (\unit{26{,}1}{kg/m^2}) e da
      estrutura de percentis por classe de IMC usada na
      Seção~\ref{sec:indice}.
\item Ittianuwat, R.; Fard, M.; Kato, K. \emph{Evaluation of Seatback
      Vibration Based on ISO 2631-1 (1997) Standard Method: The Influence of
      Vehicle Seat Structural Resonance}. Ergonomics, 2016.
\item Lee, C.T.; Daveler, B.; Candiotti, J. et al. \emph{Usability and
      Vibration Analysis of a Low-Profile Automatic Powered Wheelchair to
      Motor Vehicle Docking System}. Human Engineering Research Laboratories,
      University of Pittsburgh.
\item Muhayudin, N.A.B.; Basaruddin, K.S.; McEvoy, F.; Tansey, A.
      \emph{Development and Evaluation of a Paediatric Mid-Thoracic Spine
      Finite Element Model}. AIP Conference Proceedings, Green Design and
      Manufacture 2020, 2021.
\item Pulay, E. et al. \emph{The Effect of Additional Whole-Body Vibration on
      Musculoskeletal System in Children with Cerebral Palsy: A Systematic
      Review and Meta-Analysis of Randomized Clinical Trials}. Journal of
      Clinical Medicine, 2023.
\item Rao, S.; Sivapirakasam, S.P.; Arora, N. \emph{Whole Body Vibration
      Analysis for Children in Midi School Bus}. IOP Conference Series:
      Materials Science and Engineering, 402, 012114, 2018.
\item Rohlmann, A.; Hinz, B.; Blüthner, R.; Graichen, F.; Bergmann, G.
      \emph{Loads on a spinal implant measured in vivo during whole-body
      vibration}. European Spine Journal, 19, 1129--1135, 2010. ---
      referência [53] da ISO 2631-5:2018; única medida direta \emph{in vivo}
      revisada (Seção~\ref{sec:licoes}).
\item Schmidt, A.L.; Paskoff, G.; Shender, B.S.; Bass, C.R. \emph{Risk of
      Lumbar Spine Injury From Cyclic Compressive Loading}. Spine, 37,
      E1614--E1621, 2012. --- referência [66] da ISO 2631-5:2018; fonte
      exata da Tabela~C.1/C.2 e da Fórmula~C.4 (Seção~\ref{sec:licoes}).
\item Schust, M.; Menzel, G.; Hofmann, J.; Forta, N.G.; Pinto, I.; Hinz, B.;
      Bovenzi, M. \emph{Measures of internal lumbar load in professional
      drivers --- the use of a whole-body finite-element model for the
      evaluation of adverse health effects of multi-axis vibration}.
      Ergonomics, 58, 1191--1206, 2015. --- referência [71] da ISO
      2631-5:2018.
\item Seidel, H.; Hinz, B.; Hofmann, J.; Menzel, G. \emph{Intraspinal forces
      and health risk caused by whole-body vibration --- Predictions for
      European drivers and different field conditions}. International
      Journal of Industrial Ergonomics, 38, 856--867, 2008. --- referência
      [57] da ISO 2631-5:2018; publica a Tabela~2 usada para resolver o
      índice antropométrico (Seção~\ref{sec:indice}).
\item Seidel, H.; Pöpplau, B.M.; Morlock, M.M.; Püschel, K.; Huber, G.
      \emph{The size of lumbar vertebral endplate areas --- Prediction by
      anthropometric characteristics and significance for fatigue failure
      due to whole-body vibration}. International Journal of Industrial
      Ergonomics, 38, 844--855, 2008. --- referência [58] da ISO
      2631-5:2018; mostra que a área de placa terminal não é previsível por
      antropometria externa (Seção~\ref{sec:licoes}).
\item Sharma, A.; Kumar, A.; Mandal, B.B. \emph{Association of Risk of
      Lumbar Injury Outlined in ISO 2631-5:2018 with Self-Reported Low Back
      Pain Among Dumper Operators in Opencast Coal Mine}. Mining, Metallurgy
      \& Exploration, 2025.
\item Singh, P.; Samuel, D.V.K.; Singh, J.; Singh, R.; Prakash, C.; Dhaliwal,
      H.S. \emph{Whole Body Vibration Exposure among the Tractor Operator
      during Soil Tillage Operation: An Evaluation using ISO 2631-5
      Standard}. Applied Bionics and Biomechanics, 2022.
\item Tarabini, M.; Williams, N.; Smith, D.; Maeda, S.; Bovenzi, M.
      \emph{Scientific basis and applicability of the ISO 2631-1 health
      guidance caution zone}. International Archives of Occupational and
      Environmental Health, 99:45, 2026.
\item Zhang, Y. et al.; Rakheja, S. \emph{Modeling the Seated Human Body
      Under Whole-Body Vibration: A Multibody Dynamics Approach with Joint
      Constraints}. Multibody System Dynamics, 2026.
\item Zuska, A.; Frej, D.; Jackowski, J.; Żmuda, M. \emph{Research and
      Analysis of the Propagation of Vertical Vibrations in the Arrangement of
      a Vehicle Seat---A Child's Seat}. Sensors, 2021.
\end{enumerate}

\end{document}
